\documentclass[11pt]{article}

\usepackage[T1]{fontenc}
\usepackage[utf8]{inputenc}
\usepackage{lmodern}
\usepackage[a4paper,margin=27mm]{geometry}
\usepackage{amsmath,amssymb,amsthm,mathtools,mathrsfs}
\usepackage{microtype}
\usepackage{enumitem}
\usepackage{booktabs}
\usepackage{xcolor}
\usepackage[colorlinks=true,linkcolor=blue!45!black,citecolor=blue!45!black,
  urlcolor=blue!55!black]{hyperref}
\setlength{\emergencystretch}{2em}

\newtheorem{theorem}{Theorem}[section]
\newtheorem{proposition}[theorem]{Proposition}
\newtheorem{lemma}[theorem]{Lemma}
\newtheorem{corollary}[theorem]{Corollary}
\newtheorem*{maintheorem}{Source-modular superselection theorem}
\theoremstyle{definition}
\newtheorem{definition}[theorem]{Definition}
\theoremstyle{remark}
\newtheorem{remark}[theorem]{Remark}

\newcommand{\C}{\mathbb C}
\newcommand{\R}{\mathbb R}
\newcommand{\N}{\mathbb N}
\newcommand{\Z}{\mathbb Z}
\newcommand{\Res}{\operatorname*{Res}}
\newcommand{\ord}{\operatorname{ord}}
\newcommand{\supp}{\operatorname{supp}}
\newcommand{\Rea}{\operatorname{Re}}
\newcommand{\Ima}{\operatorname{Im}}
\newcommand{\cF}{\mathcal F}
\newcommand{\cH}{\mathcal H}
\newcommand{\cR}{\mathcal R}
\newcommand{\cS}{\mathcal S}
\newcommand{\dd}{\,\mathrm d}

\title{\textbf{The ramified-character law for the Riemann divisor}\\
\large Bifiltered source acquisition, conservative descent,
and universal boundary spectra}
\author{Michael Haufschild}
\date{}

\begin{document}

\maketitle

\begin{abstract}
We prove that a single ramified character object organizes source constructions for the Riemann
divisor which have previously appeared unrelated.  Let
\[
 \xi(s)=\tfrac12s(s-1)\pi^{-s/2}\Gamma(s/2)\zeta(s)
\]
and let \(\rho\) be a zero of multiplicity \(m_\rho\).  Attach
\[
 \chi_\rho(z)=\exp\!\bigl(z(\rho-\tfrac12)\bigr),\qquad
 \omega_{\rho,k}=
 \exp\!\left(\frac{2\pi i\,m_\rho}{2^k}\right).
\]
The first is a continuous reflection character; the second is its dyadic ramification tower.
They satisfy
\[
 \chi_\rho(z+w)=\chi_\rho(z)\chi_\rho(w),\qquad
 \chi_{1-\overline\rho}(t)=\overline{\chi_\rho(t)}^{-1},
\]
and, for \(t\ne0\),
\[
 |\chi_\rho(t)|=1\Longleftrightarrow\Rea\rho=\tfrac12,\qquad
 \omega_{\rho,k}=1\Longleftrightarrow2^k\mid m_\rho.
\]
The multiplicity-weighted infinitesimal character is not imposed on the zero set.  A singular
squarefree Möbius--Gamma source graph reconstructs
\(\mathcal G_\xi=-2\xi'/\xi\), and its balanced outer connection has residue
\[
 \Res_{s=\rho}\!\left[-(s-\tfrac12)\frac{\xi'(s)}{\xi(s)}\right]
 =-m_\rho(\rho-\tfrac12),
\]
exactly the first coefficient of the inner two-endpoint depth profile.  Positive
largest-valuation shells, compact prime-window dilation, and Hurwitz monodromy recover the same
continuous character before or independently of zero selection.

The source itself is genuinely bifiltered by valuation depth and residue-support breadth.
No additive map from this rank-two nonnegative grading to one nonnegative degree is injective;
two paired transverse scalarizations are already faithful, whereas their unpaired marginals are
not.  Tail-complete fusion across this grid forces a countably free hypertranscendental tower.
Divisor acquisition is necessarily singular as well: every nonzero nonnegative
momentum-symmetric theta-null form has dense nonclosed ordinary-\(L^2\) radical and is
nonclosable, while no entire moment symbol separates one heat node from the others.

Between source and divisor lies an odd orientation jet.  The positive Julia defect of the
Gamma-sewn prime cell has normalized normal derivative \(\log p\), while every bounded global
recovery of this prime clock has diverging norm.  Its pointed Fisher--Ward trace is odd in the
unsquared source tangent.  Squaring instead produces a canonical finite nonnegative
Xi--Hurwitz energy; this energy vanishes exactly when every \(\chi_\rho\) is unitary, but is
reflection-even and therefore detects orientation failure without choosing a side.
Squaring the completed carrier preserves its full logarithmic-current differential field; only
the odd reflection lift is lost.  A universal thermal charge law then couples that lift to
ramification: positive finite first-order density forces \(2qm=1\), so every finite-sheeted
realization selects sign monodromy and cannot descend through its full deck group.

After acquisition, every regular descendant is a shadow of \((m_\rho,\chi_\rho)\).
The real character rate is simultaneously the Wold--Poisson transport exponent, the prime-seam
Sonin/Calkin coefficient, the fixed-edge four-unit Ward obstruction, and the
finite-versus-flat Rees type.  The deck tower is simultaneously divisor-order divisibility,
analytic \(2^k\)-root descent, local loop monodromy, and the exact sharp breakdown level.
Quantitatively, the common admissible aperture set of the Poisson trace, divisor half-plane,
projective radius, and Hurwitz phase is one closed ray \([R_\xi,\infty)\).
Faithful mass deformation returns the same width with an explicit universal condition number.

We also prove an abstract conservative-descent theorem: a nonzero equivariant image of a real
exponential eigencharacter in any norm-conservative target must have zero real rate.  Hence an
off-line divisor character must either be erased or make descent nonconservative; the separate
Wold, Calkin, Ward, and Rees failures are concrete realizations of this dichotomy.
The result does not prove the Riemann hypothesis.  It identifies a complete pre-positivity
object---a bifiltered singular source, its odd orientation jet, and a ramified character
module---and proves why no scalar regular completion can replace all three stages.
\end{abstract}

\medskip
\noindent\textbf{2020 Mathematics Subject Classification.}
Primary 11M26, 11M41; Secondary 12H05, 46B99, 60E15.

\noindent\textbf{Keywords.}
Riemann hypothesis; Riemann xi function; hypertranscendence; prime valuations;
ramified character; bifiltration; conservative descent; normal cone; Rees filtration;
Hurwitz monodromy.

\section{Introduction}\label{sec:intro}

The functional equation
\[
 \xi(s)=\xi(1-s)
 \tag{1.1}\label{eq:fe}
\]
pairs a zero \(\rho=\tfrac12+c+i\gamma\) with
\(1-\rho=\tfrac12-c-i\gamma\).  The pairing preserves the divisor, but its scalar form is even:
it records the orbit and forgets which member has \(c>0\).  This loss is easy to underestimate.
Squaring an odd normal coordinate preserves the zero set.  Taking a determinant preserves
multiplicity.  A positive energy can detect that an off-line pair is present.  None of these
operations retains the sign of \(c\).

Spectral and trace-formula programs seek an operator realization from which the critical line
would follow.  Connes's trace framework and Meyer's nuclear-space realization make the
source-to-spectrum problem precise \cite{Connes1999,Meyer2005}; the scaling Hamiltonian gives a
particularly explicit arithmetic dilation picture \cite{ConnesConsani2021}.  Hedenmalm's recent
theta differential pencil isolates a related missing sesquilinear structure
\cite{Hedenmalm2026}, while Connes's recent survey emphasizes finite-to-infinite Euler
information and the Weil quadratic form \cite{Connes2026}.  The present work addresses a prior
information question:

\begin{quote}
\emph{Which source quantity still knows the transverse orientation after the completed divisor
has been formed, and which classes of transport can or cannot carry it?}
\end{quote}

The answer is not one scalar but a ramified character.  A zero of order \(m_\rho\) carries
\[
 \left(m_\rho,\chi_\rho\right),\qquad
 \chi_\rho(z)=e^{z(\rho-1/2)}.
 \tag{1.R}\label{eq:ramified-intro}
\]
The continuous character retains the full complex point and its reflection orientation.
The integer retains ramification; its dyadic deck shadows are
\(e^{2\pi i m_\rho/2^k}\).  Neither coordinate determines the other.  Their first common
coefficient is the centered logarithmic residue.

That coefficient appears first in a carrier-independent normal-cone law.
For a reflection-symmetric meromorphic carrier \(f\), set
\[
 E_\varepsilon(u)=\frac{f(\varepsilon u)}{f(u)},\qquad
 T_\varepsilon(u)=\frac{f(u)}{f((1+\varepsilon)u)}.
\]
Then \(E_\varepsilon\) exposes the reciprocal exceptional section \(f(0)/f(u)\), while
\[
 \left.\frac{\partial}{\partial\varepsilon}T_\varepsilon(u)\right|_{\varepsilon=0}
 =-u\,\frac{f'(u)}{f(u)}.
\]
Centering at the reflection fixed point yields
\[
 J_f(u)=-(u-\tfrac12)\frac{f'(u)}{f(u)}.
 \tag{1.2}\label{eq:current-intro}
\]
At a divisor point of order \(m\), its residue is the full affine vector
\(-m(\rho-\tfrac12)\).  This is the unique affine reflection-odd weight, up to scale, whose
product with the odd logarithmic derivative is reflection invariant.

The residue is also the transition coefficient between an inner and an outer chart.  Composing
literal depth steps from the two functional-equation endpoints and blowing up a moving zero
\(u=\rho+\varepsilon z\) gives
\[
 \left(\frac{z}{z+\rho-\tfrac12}\right)^m.
 \tag{1.3}\label{eq:nearby-intro}
\]
At \(z=\infty\), its first defect from the vacuum \(1\) is
\(-m(\rho-\tfrac12)z^{-1}\), exactly the residue of \(J_f\).  For \(\xi\), an actual
squarefree source graph reconstructs the outer connection \(J_\xi\) with the same normalization.

For \(\xi\), the same vector is generated from arithmetic in two independent ways.  The first is
normal to exponent depth.  Integers whose prime valuations are at most \(M\) have Dirichlet
series
\[
 A_M(s)=\frac{\zeta(s)}{\zeta((M+1)s)}.
 \tag{1.4}\label{eq:AM-intro}
\]
At fixed \(s\) this tends to \(\zeta(s)\).  At inverse depth it tends to
\(-1/(2\zeta(u))\).  Completion identifies the normal cones at \(0\) and \(1\) exactly, with
exceptional section \(1/(2\xi(u))\).

The second source is tangential to exponent depth.  Let
\[
 H_n(s)=\frac1{\zeta((n+1)s)}-\frac1{\zeta(ns)}.
 \tag{1.5}\label{eq:Hn-intro}
\]
Multiplying by \(\zeta(s)\) gives the positive Dirichlet series supported on integers with
largest prime-valuation exponent exactly \(n\).  The leading prime-two atom gives
\[
 2^{ns}H_n(s)\longrightarrow1-2^{-s},
 \qquad
 \sqrt2\,\frac{H_{n+1}(s)}{H_n(s)}
 \longrightarrow2^{1/2-s}
 \tag{1.6}\label{eq:shell-intro}
\]
for every \(\Rea s>0\).  Hence
\[
 \frac{\log|2^{1/2-s}|}{\log2}=\frac12-\Rea s.
 \tag{1.7}\label{eq:entropy-intro}
\]
This is precisely the real part of the normalized residue in
\eqref{eq:current-intro}.  It is important that \eqref{eq:shell-intro} precedes divisor
specialization: the character is forced by a positive coefficient shell.

Prime two alone aliases imaginary parts modulo \(2\pi/\log2\).  Removing its exact leading atom
from the same shell exposes prime three, producing \(3^{1/2-s}\).  Since
\(\log3/\log2\notin\mathbb Q\), the pair
\[
 s\longmapsto\bigl(2^{1/2-s},3^{1/2-s}\bigr)
 \tag{1.8}\label{eq:two-prime-intro}
\]
is injective.  The source boundary therefore retains not only horizontal displacement but the
entire complex spectral coordinate.

The third realization is local and windowed.  Dilation and inverse dilation of any
noncensoring compact bounded-variation prime window multiply the residue indexed by \(\rho\)
(at the shifted source pole \(\rho+\tfrac12\)) by
\(r^{\rho-1/2}\) and \(r^{1/2-\rho}\).  Therefore
\[
 \frac1{2\log r}\log
 \frac{|\mathcal R_{W^{\rightarrow r}}(\rho)|}
      {|\mathcal R_{W^{\leftarrow r}}(\rho)|}
 =\Rea\rho-\frac12.
 \tag{1.9}\label{eq:bv-intro}
\]
Before continuation the two windows are related by literal square-index decimation.  The
orientation is consequently source-defined, not a sign attached after locating \(\rho\).

There is also a source-side infinitesimal realization.  The Gamma-sewn scalar prime cell
\(C_x(p)\) approaches an isometry with quadratic contact, so its scalar limit loses the
first-order direction.  The positive Julia defect
\[
 D_x(p)=\sqrt{1-C_x(p)^2}
\]
restores it:
\[
 \frac{2}{\sqrt3}\frac{D_x(p)}x\longrightarrow\log p.
 \tag{1.J}\label{eq:julia-intro}
\]
On finite prime packets this is the actual logarithmic clock; globally its graph is dense in
the product Hardy space, and every bounded recovery has diverging norm.  Thus the orientation
producer is explicit but singular.  The pointed Fisher--Ward trace pairs this unsquared tangent
with double-cutoff curvature and is odd under tangent reversal.  This is the pre-divisor
operation that scalar Hellinger affinity and squared Julia coordinates erase.

The last part of the paper shows why finite-order completions cannot repair the loss.  The
geometric prime-valuation law at displacement \(c\) produces a Hellinger product.  Relative to
the Wold vacuum scale \(a_N=\prod_{p<N}(1-p^{-1})\), its two reflected faces satisfy
\[
 \begin{array}{c|c}
 c>0 & \text{flat beyond every power of }a_N,\\
 c<0 & \text{comparable to }a_N^{1/2}.
 \end{array}
 \tag{1.10}\label{eq:type-intro}
\]
The geometric law here is forced by source disintegration: it is the conditional positive-depth
fibre of the depth-one activation partition
\(A_1(\sigma)=\zeta(\sigma)/\zeta(2\sigma)\), with \(\sigma=1-2c\).
Kakutani singularity sees both faces only as zero overlap \cite{Kakutani1948}; Rees type recovers
the orientation erased by that dichotomy.  A universal operator theorem then proves that a
finite-Rees map preserves the all-orders-flat radical in any dimension.

The finite-Rees obstruction is the last member of a three-stage source-completeness law.  Before
divisor formation, exact same-prime fusion forces a countably free hypertranscendental
valuation--residue tower.  At divisor formation, the first score numerator has one less vanishing
order than its overlap, so bounded positive Gram completion is impossible and singular graph
division is forced.  After reflection, any exact sewing has essential Rees order.  These are
differential, positive, and operator shadows of one fact: the ordinary-prime source has no finite
completion compatible with fusion, divisor transgression, and internal reflection.

That statement is incomplete unless topology, source grading, ramification, and final analytic
consumer are included.  The source has a genuine depth--breadth bifiltration which no additive
scalar degree can remember.  Divisor acquisition is singular: any nonzero positive
translation-covariant form detecting the theta divisor is forced out of ordinary \(L^2\), with
dense radical and nonclosable form graph.  The acquired divisor then has two independent
coordinates.  Its continuous reflection rate is
\[
 \nu(\rho)=2(\Rea\rho-\tfrac12)
 \tag{1.11}\label{eq:rate-intro}
\]
and governs independent Wold--Poisson boundedness, compactness of the seam-weighted Sonin
corners, and positivity of the fixed-edge four-unit Ward Gram.  This coordinate is regular
exactly when \(\nu=0\).  The discrete ramification coordinate is the vanishing order
\(m_\rho\), detected by analytic root descent and local deck monodromy even when \(\nu=0\).
Theorems below prove that one scalar regular Hilbert completion cannot acquire, ramify, and
conservatively reflect this object.

\subsection*{Main result}

The paper proves the following closed theorem package.

\begin{maintheorem}\label{thm:main}
Let \(\rho\) be a right-half nontrivial zero of \(\xi\), of multiplicity \(m_\rho\), and put
\[
 c_\rho=\Rea\rho-\frac12.
\]
Then \(0<c_\rho<\tfrac12\), and:
\begin{enumerate}[label=\textup{(\roman*)}]
\item the real normal-cone residue equals \(-m_\rho c_\rho\);
\item the corrected singular squarefree source graph is exactly
\(-2\xi'/\xi\) on its Euler domain; its balanced Euler contraction is \(J_\xi\), and the
two-endpoint moving-zero profile has first outer coefficient
\(-m_\rho(\rho-\tfrac12)\);
\item the additive source-inverse/analytic-adjoint fracture is \(2c_\rho\); prime clocks
exponentiate its two branches, and the resulting canonical prime seams are not coherent;
\item the positive exponent-shell entropy equals \(-c_\rho\);
\item every noncensoring compact BV prime-window residue phase and the principal Hurwitz phase
equal \(c_\rho\);
\item the prime-two/prime-three shell spectrum recovers \(\rho\) uniquely;
\item the valuation law is the conditional geometric fibre of the depth-one squarefree member
\(A_1\) of the same exponent source;
\item its positive atoms are Fisher-damped Gamma valuation coordinates; canonical Gamma sewing
preserves prime summability exactly and its first \(N\) Rees symbols separate every \(N\)
distinct primes;
\item the positive Julia defect of the same Gamma-sewn cell is the unique
orientation-preserving two-channel lift and has normalized normal derivative \(\log p\);
on finite packets it converges to the prime logarithmic clock, while every bounded global
recovery has norm tending to infinity;
\item the pointed Fisher--Ward transgression is odd in that unsquared source tangent and is
exactly its trace pairing with double-cutoff curvature; the associated Hellinger and energy
scalarizations are even;
\item squaring the completed carrier preserves hypertranscendence and its full
logarithmic-current differential field, while the source-selected Cayley coordinate is odd,
vanishes exactly on the regular Xi divisor chart, and is invisible to every differentiable even
reader at first order; finite-dimensional jet reporters and finite local polynomial spectral
frames therefore remain incomplete;
\item the reciprocal-prime activation admits no nonzero finite recurrence; divisor
differentiation lowers overlap order by one, forces inverse-square Gram blow-up, and forbids a
uniformly bounded positive \(2\times2\) completion;
\item every tail-complete prime-residue refinement has a strict valuation-depth tower whose next
analytic layer is relatively hypertranscendental over the complete earlier prefix; exact
same-prime fusion reaches that excluded layer, so every fusion-complete differential field
contains a countably free core;
\item the reflected prime-valuation sector at \(c_\rho\) is all-orders flat, while the sector at
\(-c_\rho\) is half-normal;
\item no bounded operator family of fixed finite Rees order sews those sectors, and every exact
sewing \(R_N\) satisfies
\[
 a_N^{-K}=o(\|R_N\|)\qquad\text{for every }K\in\N.
\]
\item the Wold-depth/Poisson-time rate is the same number
\(\nu(\rho)=2c_\rho\); the two-clock energy has a finite nonzero profile on every comoving
characteristic, but is uniformly bounded on the independent nonnegative quadrant exactly when
\(\nu=0\);
\item at every positive Sonin aperture the seam-weighted two-prime Calkin corner has essential
norm
\[
 |\kappa_p(\rho)-\kappa_q(\rho)|\,(\sqrt p\sqrt q)^{-1};
\]
hence all such corners are compact exactly on the horizontal equalizer;
\item for an upper Schwarz representative of every strict-right zero, the continued Xi Ward
response violates the four-unit bound at one adapted point of the fixed Euler edge, and its
one-node boundary Gram is not positive semidefinite;
\item independently, every nonzero nonnegative momentum-symmetric Schwartz form annihilating the
Riemann theta source is an atomic real-divisor rigging with dense nonclosed ordinary-\(L^2\)
radical and is not sequentially closable through ordinary \(L^2\).
\item the divisor point carries the ramified character pair
\[
 \chi_\rho(z)=e^{z(\rho-1/2)},\qquad
 \omega_{\rho,k}=e^{2\pi i m_\rho/2^k};
\]
the first is additive, reflection sends its real-ray restriction to its conjugate inverse, and
its norm is one at one nonzero real time exactly on the critical line;
\item the absolutely convergent nonnegative Xi--Hurwitz orientation energy is the quadratic
compression of the continuous characters: it vanishes exactly when every
\(\chi_\rho\) is unitary on the real unit step, but its local summand is invariant under
critical reflection and therefore does not retain the side of a defect;
\item \(\omega_{\rho,k}^{2^k}=1\) and
\(\omega_{\rho,k}=1\) exactly when \(2^k\mid m_\rho\); the same charge is realized by every
sufficiently small descended loop and by all analytic, modulus, deck, window, and root charts at
level \(k\).  A right-half zero has persistently nontrivial rung-one loop character exactly when
its multiplicity is odd;
\item universally, a positive finite first-order thermal density for a rate-\(q\) analytic germ
forces \(2qm=1\), where \(m\) is its vanishing order.  On a finite \(r\)-sheeted branch this
selects the sign deck character and excludes full deck-invariant descent, including at the
four-unit normalization;
\item the prime source has independent valuation-depth and residue-support gradings.  No additive
scalarization \(\N^2\to\N\) is faithful, while the paired transverse degrees
\[
 (d,b)\longmapsto(d+b,d+2b)
\]
are faithful; observing the two marginal spectra without their pairing is not;
\item no entire compactified moment symbol can equal one at one positive theta heat node and zero
at every other positive heat node;
\item if
\[
 R_\xi=\sup_{\xi(\rho)=0}\left|\Rea\rho-\tfrac12\right|,
\]
then the common aperture spectrum of the Xi Poisson trace, divisor half-plane, projective radius,
and principal Hurwitz phase is exactly \([R_\xi,\infty)\);
\item a mass-\(a\) spectral deformation preserves every original aperture bound and returns it
from the deformed trace with the universal width loss
\(\sqrt{(9+a^2)/9}\);
\item universally, if \(U,V\) are real-linear actions, \(T\) intertwines them,
\(Ux=e^\nu x\), \(Tx\ne0\), and \(V\) preserves every norm, then \(\nu=0\).
\end{enumerate}
The source graph, character families, activation law, and residue-refined valuation towers in
\textup{(ii)}--\textup{(x)} are defined before the zero set is read; \(\rho\) enters only when
their boundary characters are evaluated on the divisor.
\end{maintheorem}

Theorem~\ref{thm:main} is not an RH equivalence padded with auxiliary notation.  It constructs
the missing oriented object and proves its conservation across independent source charts.  It
also proves a sharp negative result: after passage to the flat radical, neither a larger finite
matrix nor an arbitrary-dimensional bounded finite-Rees operator can recover the reflected
sector.

The later clauses give the complete phase diagram.  Wold, Calkin, Ward, and Rees failures are
horizontal shadows of one nonunitary character.  Hedenmalm's dense radical and heat-node
nonseparability belong to singular acquisition, not to that horizontal phase.  Dyadic roots and
local deck loops read multiplicity rather than displacement.  The source bifiltration explains
why one scalar completion cannot retain all stages.  Together they rule out the
assumption that one regular rank-one Hilbert completion can acquire the divisor, remember its
ramification, and orient its reflection.

\subsection*{What is and is not claimed}

The residue formula itself is classical local complex analysis.  The \(m\)-free Euler quotient,
Hellinger affinity, and Kakutani dichotomy also have classical parents
\cite{Titchmarsh1986,Hellinger1909,Kakutani1948}.  Hypertranscendence of Dirichlet families belongs
to the Ostrowski--Postnikov--Shapiro--Sparer tradition
\cite{Ostrowski1920,Postnikov1949,ShapiroSparer1986,Pong2016}.  Our contribution is the
ramified-character conservation law:

\begin{enumerate}[label=\textup{(\alph*)}]
\item one positive exponent-depth source produces the zeta bulk, reciprocal exceptional divisor,
and centered logarithmic current at three scales;
\item its squarefree score graph reconstructs the outer Xi connection, while balanced
two-endpoint depth transport supplies the inner rational normal profile, with exact transition
coefficient;
\item source inverse and analytic adjoint differ by one additive boundary character, and prime
clocks exponentiate its two branches;
\item its cofinal hypertranscendental shell character carries exactly the same signed affine
residue;
\item compact prime windows and Hurwitz monodromy recover the same character independently;
\item Fisher whitening and Gamma sewing retain the valuation source as a summability-rigid
infinite Rees resolution of unbounded exact rank;
\item the positive Julia lift restores the prime-log first jet erased by the Gamma scalar
contraction, while its global Hardy realization is necessarily singular;
\item the pointed Fisher--Ward trace is odd in that unsquared tangent and pairs it with cutoff
curvature, whereas the canonical Xi--Hurwitz energy is a faithful but reflection-even
quadratic compression;
\item evenization preserves the completed logarithmic-current field but loses the reflected
orientation lift; this parity split explains simultaneous divisor recovery and first-jet
blindness, including finite reporter and finite local-frame obstructions;
\item finite recurrence, finite differential closure, bounded positive Gram completion, and
finite-Rees reflection all fail at distinct stages of the same source;
\item reflected valuation sectors realize its sign as a finite-versus-flat asymptotic
superselection law;
\item the same charge is the Wold--Poisson transport rate and the coefficient of every
seam-weighted Sonin/Calkin corner, while the fixed-edge Ward Gram is its exact final positive
consumer;
\item positive theta-divisor detection is singular through ordinary \(L^2\), independently of
the horizontal charge;
\item the source is genuinely bifiltered: no one additive scalar degree is faithful, while two
paired transverse views recover the complete bidegree;
\item the dyadic deck tower supplies the independent ramification coordinate, explaining
analytic root descent, local loop charge, and multiplicity-sensitive Hedenmalm obstructions;
\item positive thermal boundary density couples continuous rate to discrete order by
\(2qm=1\), forcing sign monodromy and excluding full analytic deck descent;
\item the entire continuum of horizontal zero-free apertures is conserved as one closed spectrum
across divisor, Poisson, projective, and monodromy charts, with an explicit condition number
under faithful mass deformation; and
\item every faithful equivariant descent into a norm-conservative representation forces the
continuous character to be unitary.
\end{enumerate}

No assertion is made that every modular character is unitary.  Thus RH is not proved.  What is
proved is a complete classification of the information that exists before positivity: source
bidegree, odd orientation jet, divisor multiplicity, complex reflection character, global
aperture spectrum, quadratic detector, and the precise regular descents that either preserve,
erase, or fail on that object.

\section{A universal normal-cone residue law}\label{sec:normal}

This section is independent of zeta.  Let \(f\) be meromorphic on a domain stable under
\(s\mapsto1-s\), assume \(f(1-s)=f(s)\), and suppose \(f(0)\ne0\).

\begin{definition}
Define the endpoint charts, exceptional section, depth transport, and centered current by
\[
 E_\varepsilon(u)=\frac{f(\varepsilon u)}{f(u)},\qquad
 E_\varepsilon^\sharp(u)=\frac{f(1-\varepsilon u)}{f(u)},
\]
\[
 E_0^{\rm exc}(u)=\frac{f(0)}{f(u)},\qquad
 T_\varepsilon(u)=\frac{f(u)}{f((1+\varepsilon)u)},\qquad
 J_f(u)=-(u-\tfrac12)\frac{f'(u)}{f(u)}.
\]
\end{definition}

\begin{theorem}[Endpoint gluing and depth differentiation]\label{thm:endpoint}
Where the quotients are defined,
\[
 E_\varepsilon=E_\varepsilon^\sharp.
 \tag{2.1}\label{eq:endpoint-glue}
\]
If \(f\) is continuous at \(0\), then
\[
 E_\varepsilon(u)\longrightarrow E_0^{\rm exc}(u)
 \qquad(\varepsilon\to0).
 \tag{2.2}\label{eq:endpoint-limit}
\]
If \(f\) is differentiable at \(u\) and \(f(u)\ne0\), then
\[
 T_0(u)=1,\qquad
 \left.\partial_\varepsilon T_\varepsilon(u)\right|_{\varepsilon=0}
 =-u\,\frac{f'(u)}{f(u)}.
 \tag{2.3}\label{eq:depth-derivative}
\]
\end{theorem}

\begin{proof}
Equation \eqref{eq:endpoint-glue} is the functional equation applied to
\(1-\varepsilon u\).  Equation \eqref{eq:endpoint-limit} follows from
\(f(\varepsilon u)\to f(0)\).  Finally, quotient differentiation gives
\[
 \partial_\varepsilon\frac{f(u)}{f((1+\varepsilon)u)}
 =-\frac{f(u)f'((1+\varepsilon)u)u}{f((1+\varepsilon)u)^2},
\]
whose value at zero is \eqref{eq:depth-derivative}.
\end{proof}

\begin{theorem}[Affine residue conservation]\label{thm:affine-residue}
The current \(J_f\) is reflection invariant:
\[
 J_f(1-u)=J_f(u).
 \tag{2.4}\label{eq:J-even}
\]
If \(f\) has divisor order \(m\in\Z\) at \(\rho\), then
\[
 \Res_{u=\rho}J_f(u)=-m(\rho-\tfrac12).
 \tag{2.5}\label{eq:J-res}
\]
The exceptional section reverses divisor order and logarithmic connection:
\[
 \ord_\rho E_0^{\rm exc}=-m,\qquad
 \frac{(E_0^{\rm exc})'}{E_0^{\rm exc}}=-\frac{f'}f.
 \tag{2.6}\label{eq:exceptional-order}
\]
\end{theorem}

\begin{proof}
Differentiating \(f(1-u)=f(u)\) gives
\[
 \frac{f'(1-u)}{f(1-u)}=-\frac{f'(u)}{f(u)}.
\]
Multiplication by the centered affine factor proves \eqref{eq:J-even}.  Write
\(f(u)=(u-\rho)^m v(u)\), where \(v(\rho)\ne0\).  Then
\[
 \frac{f'(u)}{f(u)}=\frac{m}{u-\rho}+\frac{v'(u)}{v(u)},
\]
which proves \eqref{eq:J-res}.  Since \(f(0)\ne0\), the numerator of
\(E_0^{\rm exc}\) is a nonzero constant; \eqref{eq:exceptional-order} follows.
\end{proof}

\begin{theorem}[Inner--outer boundary-character law]\label{thm:inner-outer}
Assume that \(f\) is holomorphic near a zero \(\rho\) of multiplicity \(m\ge1\), and define the
balanced two-endpoint depth transport
\[
 B_\varepsilon^f(u)=
 \frac{f(u)}
 {f\!\left(1-(1+\varepsilon/2)
          \left(1-(1+\varepsilon/2)u\right)\right)}.
\]
For \(a=\rho-\tfrac12\) and \(z+a\ne0\),
\[
 \lim_{\varepsilon\to0,\ \varepsilon\ne0}
 B_\varepsilon^f(\rho+\varepsilon z)
 =P_{a,m}(z):=\left(\frac z{z+a}\right)^m.
 \tag{2.7}\label{eq:inner-profile}
\]
The profile and the outer logarithmic current have the same transition coefficient:
\[
 \lim_{w\to0,\ w\ne0}
 w^{-1}\!\left[P_{a,m}(w^{-1})-1\right]
 =-ma
 =\Res_{u=\rho}J_f(u).
 \tag{2.8}\label{eq:inner-outer-charge}
\]
\end{theorem}

\begin{proof}
Write \(f(u)=(u-\rho)^mg(u)\), where \(g\) is continuous at \(\rho\) and
\(g(\rho)\ne0\).  The denominator argument satisfies the exact identity
\[
 1-(1+\varepsilon/2)
       \left(1-(1+\varepsilon/2)(\rho+\varepsilon z)\right)-\rho
 =\varepsilon\!\left[
 z+a+\varepsilon(z+\rho/4)+\varepsilon^2z/4\right].
\]
Substitution into the local factorization gives \eqref{eq:inner-profile}.  In the inverse
coordinate,
\[
 P_{a,m}(w^{-1})=(1+aw)^{-m}=1-maw+O(w^2),
\]
which proves the first equality in \eqref{eq:inner-outer-charge}; the second is
Theorem~\ref{thm:affine-residue}.
\end{proof}

\begin{theorem}[Additive--multiplicative duality--adjunction fracture]
\label{thm:duality-fracture}
The profile with shift \(b\) has pole at \(-b\).  Source inversion applies \(a\mapsto-a\),
whereas analytic adjunction applies \(a\mapsto\overline a\); their resulting nearby poles are
\[
 \pi_{\rm inv}(a)=a,\qquad \pi_{\rm adj}(a)=-\overline a.
\]
Their additive fracture is
\[
 \Delta(\rho):=\pi_{\rm inv}(a)-\pi_{\rm adj}(a)
 =a+\overline a=2(\Rea\rho-\tfrac12).
 \tag{2.9}\label{eq:additive-fracture}
\]
Consequently
\[
 \Delta(\rho)=0\Longleftrightarrow\Rea\rho=\tfrac12,\qquad
 \Rea\Res_\rho J_f=-\frac m2\,\Delta(\rho).
 \tag{2.10}\label{eq:fracture-residue}
\]
For every prime \(p\), the corresponding multiplicative fracture is the exponential image of
the same two poles:
\[
 D_p(\rho)
 =p^{\pi_{\rm inv}(a)}-p^{\pi_{\rm adj}(a)}
 =p^{\rho-1/2}-p^{1/2-\overline\rho}.
 \tag{2.11}\label{eq:prime-fracture}
\]
Put
\[
 \kappa_p(\rho)=-(\log p)p^{-1/2}D_p(\rho).
 \]
If \(0<\Rea\rho<1\), then the family \((\kappa_p(\rho))_p\) is constant over the primes if and
only if \(\Rea\rho=\tfrac12\); in that case every member is zero.
\end{theorem}

\begin{proof}
Equations \eqref{eq:additive-fracture}--\eqref{eq:prime-fracture} are direct.  If
\(\Rea\rho=\tfrac12\), then \(a=-\overline a\), so every fracture vanishes.  Conversely, assume
all \(\kappa_p(\rho)\) have one value.  The function
\[
 \kappa_x(\rho)
 =(\log x)\left(x^{-\overline\rho}-x^{-(1-\rho)}\right)
\]
tends to zero as \(x\to+\infty\) throughout the open strip.  Arbitrarily large primes force the
common value to be zero.  Its \(p=2\) member vanishes only when the two positive-real
exponentials have equal moduli, namely when
\(2^{-\Rea\rho}=2^{-(1-\Rea\rho)}\).  Hence \(\Rea\rho=\tfrac12\).
\end{proof}

\begin{proposition}[Uniqueness of centered affine orientation]\label{prop:unique-affine}
Let \(w(u)=au+b\).  Then
\[
 w(1-u)=-w(u)\quad\text{for all }u
\]
if and only if \(b=-a/2\).  Hence every affine reflection-odd weight is a scalar multiple of
\(u-\tfrac12\).
\end{proposition}

\begin{proof}
Substitute \(u=\tfrac12\) to obtain \(a/2+b=0\).  The converse is immediate.
\end{proof}

\begin{remark}
The current \(J_f\) is even because it is a product of two odd factors.  Its residue is not even
data: the affine vector in \eqref{eq:J-res} changes sign under reflection.  Squaring that vector
retains divisor location but erases the only affine orientation allowed by
Proposition~\ref{prop:unique-affine}.
\end{remark}

\begin{remark}[Why the balanced scalar is not a Herglotz carrier]
In a centered real coordinate, a reflection-even entire carrier with a nonzero real zero
\(a>0\) has a paired zero at \(-a\).  The Euler-contracted connection has residues \(-ma\) and
\(+ma\).  A meromorphic Herglotz coefficient requires one residue sign at every real pole, and an
anti-Herglotz coefficient requires the opposite sign.  Thus neither convention can contain both
residues.  The obstruction disappears for the uncontracted one-form \(-d\log f\), whose residues
all equal \(-m\).  Theorem~\ref{thm:duality-fracture} explains the sign reversal: Euler
contraction retains the additive fracture rather than making it positive.
\end{remark}


\section{Positive exponent depth and the reciprocal zeta boundary}\label{sec:depth}

For \(M\ge0\), let \(b_M(n)\) be the indicator that every prime exponent of \(n\) is at most
\(M\).  Equivalently, \(n\) is \((M+1)\)-free.

\begin{proposition}[Bounded-exponent Euler source]\label{prop:bounded-euler}
For \(\Rea s>1\),
\[
 A_M(s):=\sum_{n\ge1}\frac{b_M(n)}{n^s}
 =\prod_p\frac{1-p^{-(M+1)s}}{1-p^{-s}}
 =\frac{\zeta(s)}{\zeta((M+1)s)}.
 \tag{3.1}\label{eq:bounded-euler}
\]
\end{proposition}

\begin{proof}
At each prime the allowed valuations are \(0,\dots,M\), so the local sum is
\[
 1+p^{-s}+\cdots+p^{-Ms}
 =\frac{1-p^{-(M+1)s}}{1-p^{-s}}.
\]
Absolute convergence permits multiplication over primes.
\end{proof}

\begin{theorem}[Three-scale source law]\label{thm:three-scale}
The same positive source has the following limits.
\begin{enumerate}[label=\textup{(\roman*)}]
\item If \(\Rea s>0\) and \(s\ne1\), then, for the meromorphic continuation of
\(A_M\),
\[
 A_M(s)\longrightarrow\zeta(s).
 \tag{3.2}\label{eq:fixed-limit}
\]
\item For every \(u\in\C\), in the meromorphic sense,
\[
 A_M\!\left(\frac{u}{M+1}\right)
 =\frac{\zeta(u/(M+1))}{\zeta(u)}
 \longrightarrow-\frac1{2\zeta(u)}.
 \tag{3.3}\label{eq:normal-limit}
\]
\item If the observation scale is \(u/(M+1)^2\), then
\[
 A_M\!\left(\frac{u}{(M+1)^2}\right)\longrightarrow1.
 \tag{3.4}\label{eq:micro-limit}
\]
\end{enumerate}
\end{theorem}

\begin{proof}
For \eqref{eq:fixed-limit}, \(\Rea((M+1)s)\to+\infty\), so
\(\zeta((M+1)s)\to1\).  Equation \eqref{eq:normal-limit} is exact before taking the
limit; use \(\zeta(0)=-1/2\).  In \eqref{eq:micro-limit}, both the numerator argument and the
denominator argument tend to zero, and their limiting values are the same nonzero number
\(-1/2\).
\end{proof}

\begin{corollary}[Completed endpoint normal cone]\label{cor:completed-normal}
Put
\[
 X_\varepsilon(u)=\frac{\xi(\varepsilon u)}{\xi(u)},\qquad
 X_\varepsilon^\sharp(u)=\frac{\xi(1-\varepsilon u)}{\xi(u)}.
\]
Then
\[
 X_\varepsilon=X_\varepsilon^\sharp,\qquad
 X_\varepsilon(u)\longrightarrow\frac1{2\xi(u)}.
 \tag{3.5}\label{eq:completed-normal}
\]
At a zero \(\rho\) of multiplicity \(m_\rho\), the exceptional section has a pole of order
\(m_\rho\).
\end{corollary}

\begin{proof}
Apply Theorems~\ref{thm:endpoint} and \ref{thm:affine-residue} to \(f=\xi\), using
\(\xi(0)=\xi(1)=1/2\).
\end{proof}

\begin{corollary}[Source-generated Xi current]\label{cor:xi-current}
Consecutive completed depths generate the Xi logarithmic current:
\[
 \left.\partial_\varepsilon
 \frac{\xi(u)}{\xi((1+\varepsilon)u)}\right|_{\varepsilon=0}
 =-u\,\frac{\xi'(u)}{\xi(u)}.
 \tag{3.6}\label{eq:xi-depth-current}
\]
The reflection-balanced current
\[
 J_\xi(u)=-(u-\tfrac12)\frac{\xi'(u)}{\xi(u)}
 \]
has residue \(-m_\rho(\rho-\tfrac12)\) at every zero \(\rho\).
\end{corollary}

The point of Corollary~\ref{cor:xi-current} is structural.  The determinant and its logarithmic
connection are not introduced after enumerating zeros.  They are the exceptional state and
infinitesimal depth generator of the literal bounded-exponent source \eqref{eq:bounded-euler}.

\section{Positive hypertranscendental shells}\label{sec:shells}

Let
\[
 M(n)=\max_p v_p(n),\qquad n\ge2,
\]
and put \(M(1)=0\).  Write \(a_m(n)=1_{\{M(n)=m\}}\).  Then the \(a_m\) are pairwise disjoint
\(0/1\)-valued sources and
\[
 \sum_{m\ge1}a_m(n)=1\qquad(n\ge2).
 \tag{4.1}\label{eq:shell-partition}
\]

\begin{proposition}[Shell factorization]\label{prop:shell-factor}
For \(\Rea s>1\),
\[
 S_m(s):=\sum_{n\ge1}\frac{a_m(n)}{n^s}
 =\zeta(s)H_m(s),
 \qquad
 H_m(s)=\frac1{\zeta((m+1)s)}-\frac1{\zeta(ms)}.
 \tag{4.2}\label{eq:shell-factor}
\]
For every \(m\ge12\), \(H_m\) is holomorphic and zero-free on
\(\Rea s>\tfrac12\).
\end{proposition}

\begin{proof}
The identity \(a_m=b_m-b_{m-1}\), with the indexing translated to
Proposition~\ref{prop:bounded-euler}, gives \eqref{eq:shell-factor}.  It remains to prove
nonvanishing.  Write \(s=\sigma+it\), \(x=m\sigma>6\), and
\[
 \zeta(ms)-\zeta((m+1)s)=A+T,
\]
\[
 A=2^{-ms}-2^{-(m+1)s},\qquad
 T=\sum_{n\ge3}\bigl(n^{-ms}-n^{-(m+1)s}\bigr).
\]
Since \(2^{-\sigma}<2^{-1/2}<71/100\),
\[
 |A|=2^{-m\sigma}|1-2^{-s}|
 >\frac{29}{100}\,2^{-m\sigma}.
\]
For \(n\ge3\),
\[
 |n^{-ms}-n^{-(m+1)s}|
 \le n^{-x}(1+n^{-\sigma})<\frac85n^{-x}
 \le\frac85\,2^{-x}\frac{64}{n^6}.
\]
The integral estimate
\[
 \sum_{n\ge3}\frac1{n^6}
 \le\frac1{3^6}+\frac1{4^6}+\frac1{2500}
\]
therefore gives
\[
 |T|<\frac{21}{100}\,2^{-m\sigma}<|A|.
\]
Thus \(\zeta(ms)\ne\zeta((m+1)s)\).  Both zeta factors are holomorphic and nonzero because
their arguments have real part greater than one, and
\[
 H_m(s)=
 \frac{\zeta(ms)-\zeta((m+1)s)}
      {\zeta(ms)\zeta((m+1)s)}.
\]
\end{proof}

\begin{remark}
The threshold \(12\) is an explicit uniform threshold, not an optimality claim.
\end{remark}

\subsection{Differential freeness}

We give the algebraic argument because ordinary independence of the shells would not suffice:
the common zeta factor must be removed before significance is measured.

Let \(\mathcal A\) be the \(\C\)-algebra of arithmetic functions under Dirichlet convolution.
It is an integral domain.  Write
\[
 \Omega(n)=\sum_pv_p(n).
\]
For a prime \(q\), define
\[
 (\partial_q f)(n)=v_q(nq)f(nq),\qquad
 (L^kf)(n)=(-\log n)^k f(n).
 \tag{4.3}\label{eq:prime-derivation}
\]
A coefficient check gives
\(\partial_q(f*g)=(\partial_qf)*g+f*(\partial_qg)\), while Dirichlet transformation sends
\(L^k\) to the \(k\)-th analytic derivative.

\begin{theorem}[\(\Omega\)-leading-echelon criterion]\label{thm:omega}
Let \(R\) be finite and nonempty, let \(\ell:R\to\N_{>0}\) be injective, and let
\(f_r\in\mathcal A\).  Suppose that for every \(r\in R\), prime \(q\), and \(k,n\ge0\),
\begin{align}
 \Omega(n)<\ell(r)-1
 &\Longrightarrow \partial_q(L^kf_r)(n)=0,
 \tag{4.4}\label{eq:omega-below}\\
 \Omega(n)=\ell(r)-1
 &\Longrightarrow
 \partial_q(L^kf_r)(n)
 =\ell(r)(-\ell(r)\log q)^k
 \mathbf1_{\{q^{\ell(r)-1}\}}(n).
 \tag{4.5}\label{eq:omega-leading}
\end{align}
Then, for every \(J\ge0\), the rectangle
\[
 \{L^kf_r:r\in R,\ 0\le k\le J\}
 \tag{4.6}\label{eq:omega-rectangle}
\]
is algebraically independent over \(\C\) in \(\mathcal A\).  If all \(f_r\) have a finite
half-plane of absolute Dirichlet convergence, their transforms are jointly differentially
algebraically independent over \(\C(s)\).
\end{theorem}

\begin{proof}
Index the columns by \(I=R\times\{0,\ldots,J\}\).  For every \(i=(r,t)\), choose a distinct
prime \(q_i\) and form the \(\mathcal A\)-valued Jacobian
\[
 \mathcal J_{i,(u,k)}=\partial_{q_i}(L^kf_u).
\]
Put \(e(u,k)=\ell(u)-1\) and
\[
 N=\prod_{i=(r,t)\in I}q_i^{e(r,t)}.
\]
Every convolution term contributing to \(\det(\mathcal J)(N)\) has total
\(\Omega\)-grade \(\sum_i e(i)\).  By \eqref{eq:omega-below}, no factor lies below its column
grade.  Consequently every factor lies in its lowest possible grade, where
\eqref{eq:omega-leading} computes it exactly.

Unique factorization forces each contributing permutation to preserve equal-\(e\) blocks.
Injectivity of \(\ell\) identifies these with the blocks belonging to one \(u\in R\).
Up to a nonzero scalar, the \(u\)-block is
\[
 \left[(-\ell(u)\log q_{(u,t)})^k\right]_{0\le t,k\le J}.
\]
Its determinant is the nonzero Vandermonde product
\[
 (-\ell(u))^{J(J+1)/2}
 \prod_{0\le t<v\le J}
 \bigl(\log q_{(u,v)}-\log q_{(u,t)}\bigr).
\]
Thus \(\det\mathcal J\ne0\) in \(\mathcal A\).  The characteristic-zero Jacobian criterion
implies algebraic independence: a polynomial relation of least total degree would put all of
its nonzero partial derivatives in the kernel of an invertible Jacobian over
\(\operatorname{Frac}(\mathcal A)\).

It remains to justify the coefficient field \(\C(s)\).  After clearing denominators, a relation
over \(\C(s)\) has the form
\[
 \sum_{n\ge1}Q_n(s)n^{-s}=0,
 \tag{4.7}\label{eq:polynomial-dirichlet}
\]
where the degrees of the polynomials \(Q_n\) are uniformly bounded.  Choose the least \(n\)
with \(Q_n\ne0\), multiply by \(n^s\), and let \(s=\sigma\to+\infty\).
Exponential separation dominates fixed polynomial growth, forcing \(Q_n=0\), a contradiction.
Comparing powers of \(s\) then returns a relation in \(\mathcal A\), already excluded.
\end{proof}

\begin{remark}
Only the lowest nonzero \(\Omega\)-layer enters the proof.  Arbitrary higher support can be
added without altering the conclusion.  This upper-tail immunity is the reason the theorem
applies to complete exponent shells rather than isolated prime powers.
\end{remark}

\begin{lemma}[Prime-ray leading layer]\label{lem:prime-ray}
If \(a_m(n)\ne0\), then \(\Omega(n)\ge m\), and equality forces
\(n=p^m\) for a unique prime \(p\).  Consequently, for every prime \(q\),
\[
 \Omega(n)<m-1\Longrightarrow\partial_q(L^ka_m)(n)=0,
\]
whereas at grade \(m-1\),
\[
 \partial_q(L^ka_m)(n)
 =m(-m\log q)^k\mathbf1_{\{q^{m-1}\}}(n).
 \tag{4.8}\label{eq:prime-ray}
\]
\end{lemma}

\begin{proof}
If \(a_m(n)\ne0\), some \(p^m\) divides \(n\), giving \(\Omega(n)\ge m\).
Equality in \(n=p^mu\) forces \(\Omega(u)=0\), hence \(u=1\).  If
\(\partial_q(L^ka_m)(n)\ne0\), then \(a_m(nq)\ne0\), so
\(m\le\Omega(n)+1\).  At equality, \(nq=p^m\); divisibility forces \(p=q\) and
\(n=q^{m-1}\), where direct substitution gives \eqref{eq:prime-ray}.
\end{proof}

\begin{theorem}[Relative shell hypertranscendence]\label{thm:relative-ht}
Let
\[
 K_\zeta=\C(s)\bigl(\zeta,\zeta',\zeta'',\ldots\bigr).
\]
The countable family \((H_m)_{m\ge2}\) is differentially algebraically independent over
\(K_\zeta\).  Equivalently, no nonzero polynomial differential relation involving finitely many
shell multipliers survives after rational functions of \(s\), \(\zeta\), and all derivatives of
\(\zeta\) are admitted as coefficients.
\end{theorem}

\begin{proof}
Take any finite family and derivative rectangle.  Assign level \(1\) to the constant coefficient
source of \(\zeta\), and level \(m\) to \(a_m\).  The level-one source satisfies
\[
 \partial_q(L^k\mathbf1)(n)
 =(-\log q)^k\mathbf1_{\{1\}}(n)
 \quad\text{when }\Omega(n)=0,
\]
and Lemma~\ref{lem:prime-ray} handles every shell level.  Theorem~\ref{thm:omega} proves joint
differential independence of \(\zeta\) and the \(S_m\) over \(\C(s)\).

By the tower property, the \(S_m\) remain jointly differentially independent after adjoining the
complete differential field generated by \(\zeta\).  Multiplication by
\(\zeta^{-1}\in K_\zeta^\times\) is an invertible triangular change of jet coordinates:
\[
 (\zeta^{-1}S_m)^{(k)}
 =\zeta^{-1}S_m^{(k)}
  +\sum_{j<k}\binom{k}{j}(\zeta^{-1})^{(k-j)}S_m^{(j)}.
\]
Since \(\zeta^{-1}S_m=H_m\) as meromorphic germs, the \(H_m\) are jointly differentially
independent over \(K_\zeta\).  Finite character proves the countable assertion.  This is a
Dirichlet-series leading-echelon argument in the Ostrowski--Shapiro--Sparer tradition
\cite{Ostrowski1920,ShapiroSparer1986}.
\end{proof}

\subsection{The source is bifiltered, not deeply filtered}

The shell index records valuation depth, but it is not the only source grade.  Fix finitely many
residue coordinates and write \(S\) for the support on which a pointed residue observable is
nontrivial.  A prime-local carrier then has two independent labels:
\[
 (e,S),\qquad e\in\N,\quad S\subseteq I.
\]
For carriers \(F_{e,S}\) and \(F_{r,T}\), exact same-prime fusion satisfies
\[
 F_{e,S}\star F_{r,T}=F_{e+r+1,S\cup T}.
 \tag{4.B1}\label{eq:bifilter-fusion}
\]
On disjoint pure supports, shifted depth \(d=e+1\) and breadth \(b=|S|\) are both additive.
If \(S\cap T\ne\varnothing\), depth still adds but breadth loses principal degree:
\[
 |S\cup T|<|S|+|T|.
 \tag{4.B2}\label{eq:breadth-drop}
\]
This is the source-level reason that shell depth and prime-normal codimension cannot be identified.

\begin{theorem}[Bifiltered scalarization law]\label{thm:bifilter}
No additive map \(w:\N^2\to\N\) is injective.  In contrast, the paired transverse map
\[
 \Phi(d,b)=(d+b,d+2b)
 \tag{4.B3}\label{eq:paired-degree}
\]
is injective.  Pairing is essential: the two packets
\[
 \{(1,1),(2,1)\},\qquad \{(0,2),(3,0)\}
\]
have the same ordinary-total-degree multiset and the same \(d+2b\) multiset, but with opposite
cross-pairing.
\end{theorem}

\begin{proof}
Put \(a=w(1,0)\) and \(b=w(0,1)\).  If \(a=0\) or \(b=0\), one coordinate axis already collides
with the origin.  Otherwise additivity gives
\[
 w(b,0)=ba=ab=w(0,a),
\]
although \((b,0)\ne(0,a)\).  Thus \(w\) is never injective.  If
\(\Phi(d,b)=\Phi(d',b')\), subtracting the first coordinate equality from the second gives
\(b=b'\), and then \(d=d'\).  The final assertion follows by direct evaluation:
ordinary degrees are \(\{2,3\}\) for both packets and transverse degrees are \(\{3,4\}\) for
both, with the indicated swapped pairing.
\end{proof}

Theorem~\ref{thm:bifilter} changes the interpretation of
Theorem~\ref{thm:relative-ht}.  Hypertranscendental escape is not merely ``more depth''.  Fusion
moves in a depth--breadth grid, and no scalar filtration can be faithful to that motion.
The positive replacement is exactly two paired transverse realizations.  A scalar conormal can
therefore carry the divisor coefficient, but it cannot by itself generate the independent
prime-labelled normal frame required by a Sonin/Calkin consumer.  That frame is not extra
decoration; it is the second source grade.

\subsection{Least-atom characters}

\begin{lemma}[Prime-two extraction]\label{lem:prime-two}
If \(\Rea w>1\), then
\[
 \left|2^w(\zeta(w)-1)-1\right|
 \le\frac2{\Rea w-1}.
 \tag{4.9}\label{eq:two-bound}
\]
\end{lemma}

\begin{proof}
Separate \(2^{-w}\) from the Dirichlet series and bound the remaining sum by the integral test:
\[
 \left|\sum_{n\ge3}n^{-w}\right|
 \le\sum_{n\ge3}n^{-\Rea w}
 \le\frac{2^{1-\Rea w}}{\Rea w-1}.
\]
Multiplication by \(|2^w|=2^{\Rea w}\) gives \eqref{eq:two-bound}.
\end{proof}

\begin{theorem}[Shell-transfer character]\label{thm:shell-character}
For every \(s\) with \(\Rea s>0\),
\[
 2^{ns}H_n(s)\longrightarrow1-2^{-s},
 \tag{4.10}\label{eq:first-atom}
\]
\[
 \frac{H_{n+1}(s)}{H_n(s)}\longrightarrow2^{-s},
 \qquad
 \sqrt2\,\frac{H_{n+1}(s)}{H_n(s)}
 \longrightarrow\chi_2(s):=2^{1/2-s}.
 \tag{4.11}\label{eq:adjacent-shell}
\]
Moreover,
\[
 \frac{\log|\chi_2(s)|}{\log2}=\frac12-\Rea s,
 \qquad
 |\chi_2(s)|=1\iff\Rea s=\frac12.
 \tag{4.12}\label{eq:shell-entropy}
\]
\end{theorem}

\begin{proof}
Lemma~\ref{lem:prime-two}, applied to \(w=ns\), gives
\[
 2^{ns}\bigl(\zeta(ns)-1\bigr)\to1.
\]
Since \(\zeta(ns)^{-1}\to1\),
\[
 2^{ns}\bigl(\zeta(ns)^{-1}-1\bigr)\to-1.
\]
Taking the adjacent difference in \(n\) yields \eqref{eq:first-atom}.  The limit
\(1-2^{-s}\) is nonzero for \(\Rea s>0\), so the ratio of consecutive normalized
expressions gives \eqref{eq:adjacent-shell}.  Finally
\[
 |\chi_2(s)|=2^{1/2-\Rea s},
\]
which proves \eqref{eq:shell-entropy}.
\end{proof}

Remove the exact prime-two asymptotic atom and repeat the argument with prime three.  Define
\[
 \widetilde H_n(s)
 =H_n(s)-2^{-ns}(1-2^{-s}).
\]

\begin{theorem}[Two-prime shell tomography]\label{thm:two-prime}
For \(\Rea s>0\),
\[
 3^{ns}\widetilde H_n(s)\longrightarrow1-3^{-s},
 \qquad
 \sqrt3\,\frac{\widetilde H_{n+1}(s)}{\widetilde H_n(s)}
 \longrightarrow\chi_3(s):=3^{1/2-s}.
 \tag{4.13}\label{eq:three-character}
\]
The map
\[
 s\longmapsto(\chi_2(s),\chi_3(s))
 \tag{4.14}\label{eq:two-prime-map}
\]
is injective on \(\C\).
\end{theorem}

\begin{proof}
The first two statements follow from the same tail estimate after extracting the \(2^{-w}\) and
\(3^{-w}\) terms of \(\zeta(w)^{-1}\).  If the two character pairs agree at \(s\) and \(t\),
their norms give \(\Rea s=\Rea t\).  Their phases give
\[
 (\Ima s-\Ima t)\log2\in2\pi\Z,\qquad
 (\Ima s-\Ima t)\log3\in2\pi\Z.
\]
If \(\Ima s\ne\Ima t\), division would imply \(\log3/\log2\in\mathbb Q\), hence
\(3^a=2^b\) for nonzero integers \(a,b\), impossible by unique factorization.  Thus
\(\Ima s=\Ima t\).
\end{proof}

\section{Transgression to the Xi divisor}\label{sec:transgression}

Apply Section~\ref{sec:normal} to \(f=\xi\).

\begin{theorem}[Normal-to-cofinal modular transgression]\label{thm:transgression}
Let \(\rho\) be a zero of \(\xi\) of multiplicity \(m_\rho\).  Then
\[
 \Rea\Res_{s=\rho}J_\xi(s)
 =m_\rho\,\frac{\log|\chi_q(\rho)|}{\log q}
 =m_\rho\left(\frac12-\Rea\rho\right)
 \tag{5.1}\label{eq:transgression}
\]
for every \(q>1\), where \(\chi_q(s)=q^{1/2-s}\).
For \(m_\rho\ne0\), the following are equivalent:
\[
 \Rea\Res_{\rho}J_\xi=0,\qquad
 |\chi_q(\rho)|=1,\qquad
 \Rea\rho=\frac12.
 \tag{5.2}\label{eq:neutral}
\]
\end{theorem}

\begin{proof}
Theorem~\ref{thm:affine-residue} gives
\[
 \Rea\Res_\rho J_\xi=-m_\rho(\Rea\rho-\tfrac12).
\]
For \(q>1\),
\[
 \log|q^{1/2-\rho}|=(\tfrac12-\Rea\rho)\log q.
\]
This proves \eqref{eq:transgression} and \eqref{eq:neutral}.
\end{proof}

\begin{remark}
Equation \eqref{eq:transgression} is independent of the atom base.  What is nontrivial in the
arithmetic realization is that \(q=2\) and \(q=3\) are not freely chosen coordinates: they are
the first two cofinal atoms extracted from the same positive hypertranscendental shell source.
\end{remark}

\subsection{Compact prime-window recovery}

Let \(W:\R\to\C\) have bounded variation and compact support in
\([a,b]\subset(0,\infty)\), and put
\[
 I_W(s)=\int_0^\infty W(x)x^{-s}\,\dd x.
 \tag{5.3}\label{eq:window-mellin}
\]
The associated moving-prime and prime--polar coefficients are
\[
 P_W(n)=\sum_{k\ge1}\frac{\Lambda(k)}{\sqrt{k}}W(n/k),\qquad
 C_W(n)=2I_W(\tfrac32)\sqrt n-2P_W(n).
 \tag{5.4}\label{eq:window-coefficients}
\]
The Dirichlet series of \(C_W\) is absolutely convergent for \(\Rea s>3/2\) and has a canonical
meromorphic continuation.  At a nontrivial zeta zero \(\rho\) of multiplicity \(m_\rho\), its
Laurent residue at \(\rho+\tfrac12\) is
\[
 \mathcal R_W(\rho)=2m_\rho I_W(\rho+\tfrac12).
 \tag{5.5}\label{eq:window-base-residue}
\]
When this value is nonzero, the continued source has a simple pole there; we then call \(W\)
noncensoring at \(\rho\).  For \(r>1\), define
\[
 W^{\rightarrow r}(x)=W(rx),\qquad
 W^{\leftarrow r}(x)=W(x/r).
\]

\begin{theorem}[Window-independent residue phase]\label{thm:bv-phase}
The continued residues satisfy
\[
 \mathcal R_{W^{\rightarrow r}}(\rho)
 =r^{\rho-1/2}\mathcal R_W(\rho),\qquad
 \mathcal R_{W^{\leftarrow r}}(\rho)
 =r^{1/2-\rho}\mathcal R_W(\rho).
 \tag{5.6}\label{eq:window-residues}
\]
Consequently
\[
 \frac1{2\log r}\log
 \frac{|\mathcal R_{W^{\rightarrow r}}(\rho)|}
      {|\mathcal R_{W^{\leftarrow r}}(\rho)|}
 =\Rea\rho-\frac12.
 \tag{5.7}\label{eq:window-phase}
\]
The value is independent of \(W\) and \(r\).
\end{theorem}

\begin{proof}
The substitutions \(y=rx\) and \(y=x/r\) give
\[
 I_{W^{\rightarrow r}}(s)=r^{s-1}I_W(s),\qquad
 I_{W^{\leftarrow r}}(s)=r^{1-s}I_W(s).
\]
Evaluate at \(s=\rho+\tfrac12\) in \eqref{eq:window-base-residue}.  This gives
\eqref{eq:window-residues}; taking norms and logarithms gives \eqref{eq:window-phase}.
\end{proof}

For integer \(r>1\), both coefficient streams satisfy square-index decimation before
continuation:
\[
 P_{W^{\rightarrow r}}(n)=P_{W^{\leftarrow r}}(r^2n),\qquad
 C_{W^{\rightarrow r}}(n)=C_{W^{\leftarrow r}}(r^2n).
 \tag{5.8}\label{eq:square-decimation}
\]
Thus \eqref{eq:window-phase} is not a divisor-defined normalization.

\subsection{Hurwitz monodromy}

The centered complex dilation character at \(\rho\) is
\[
 \chi_\rho(z)=e^{z(\rho-1/2)}.
 \]
Its real ray gives the Mellin norm growth, while its imaginary period gives
\[
 \chi_\rho(-2\pi i)
 =-e^{-2\pi i\rho}.
 \tag{5.9}\label{eq:hurwitz}
\]
For \(0<\Rea\rho<1\), the principal argument has no aliasing in the transverse coordinate:
\[
 \frac{|\arg\chi_\rho(-2\pi i)|}{2\pi}
 =\left|\Rea\rho-\frac12\right|.
 \tag{5.10}\label{eq:hurwitz-phase}
\]
For a right-half zero, \eqref{eq:window-phase}, \eqref{eq:hurwitz-phase}, and the negative of
the shell entropy \eqref{eq:shell-entropy} all equal the same positive number
\[
 c_\rho=\Rea\rho-\frac12.
 \tag{5.11}\label{eq:charge}
\]

\section{Prime-valuation type fracture}\label{sec:fracture}

The previous sections identify the charge.  We now show that its sign is an asymptotic type, not
a finite grading choice.  The valuation sector is not an auxiliary model: it is the conditional
depth fibre of the same Euler source.

For \(\sigma>1\), diagonal activation has the squarefree partition
\[
 \mathcal Z(\sigma)
 =A_1(\sigma)=\frac{\zeta(\sigma)}{\zeta(2\sigma)}
 =\prod_p(1+p^{-\sigma}).
 \tag{6.1}\label{eq:activation}
\]
Write \(q=p^{-\sigma}\).  Normalizing one local factor gives vacant/occupied masses
\(1/(1+q)\) and \(q/(1+q)\).  Disintegrating the occupied state by its geometric valuation depth
gives
\[
 \nu_q(0)=\frac1{1+q},\qquad
 \nu_q(k)=\frac{(1-q)q^k}{1+q}\quad(k\ge1).
 \tag{6.2}\label{eq:activation-law}
\]
Its occupancy probability and conditional positive-depth law are
\[
 \mathbb P(k\ge1)=\frac q{1+q},\qquad
 \mathbb P(k=e+1\mid k\ge1)=(1-q)q^e.
 \tag{6.3}\label{eq:conditional-depth}
\]
Thus squarefree occupancy is the Bernoulli base and valuation depth is its geometric fibre.
This lies in the classical probabilistic Euler-product setting \cite{ArratiaBarbourTavare2003};
the use below is the reflected Rees classification of that exact fibre.
The occupancy Fisher variance is \(q/(1+q)^2\); at a prime, the unnormalized positive-depth atoms
are exactly the Möbius shell differences.  Identities \eqref{eq:activation-law}--%
\eqref{eq:conditional-depth} remain literal finite-prime probability laws for every
\(0<q<1\), whether or not the global Euler product converges.

\subsection{Fisher whitening is an infinite Rees resolution}

At the critical node \(q_p=p^{-1}\), put
\[
 b_{p,e}=\frac{(1-q_p)q_p^{e+1}}{1+q_p},\qquad
 \delta_p=\frac{\sqrt{q_p}}{1+q_p},\qquad
 g_{p,e}=\sqrt{q_p}(1-q_p)q_p^e.
 \tag{6.4}\label{eq:fisher-gamma}
\]
Thus \(b_{p,e}\) is the positive atom in \eqref{eq:activation-law},
\(\delta_p^2=q_p/(1+q_p)^2\) is the occupancy Fisher variance, and \(g_{p,e}\) is the
half-shifted Gamma valuation coordinate.

For aperture \(x>0\), set \(q_{x,p}=p^{-x}\) and
\(\phi_\infty(t)=\cosh(t)^{-1/2}\).  The source-normalized doubled-frequency Gamma multiplier is
\[
 C_x(p)=\frac12\,\phi_\infty(2x\log p)^{1/2}
             \frac{1+q_{x,p}}{\sqrt{q_{x,p}}}.
 \tag{6.5}\label{eq:gamma-multiplier}
\]
Define the sewn valuation ray
\[
 F_{x,p}(e)=C_x(p)\frac{q_p}{1+q_p}(1-q_{x,p})q_{x,p}^e,
 \qquad e\in\N.
 \tag{6.6}\label{eq:gamma-ray}
\]
Write \(\Delta f(e)=f(e+1)-f(e)\).

\begin{theorem}[Fisher--Gamma Rees resolution]\label{thm:gamma-rees}
The activation and Gamma coordinates obey:
\begin{enumerate}[label=\textup{(\roman*)}]
\item \(b_{p,e}=\delta_p g_{p,e}\), and
\[
 \sum_{e\ge0}g_{p,e}^2
 =\frac{q_p(1-q_p)}{1+q_p}
 =\frac{p-1}{p(p+1)}.
 \tag{6.7}\label{eq:gamma-mass}
\]
Consequently the Fisher-whitened fibre has divergent total mass over the primes.
\item The exact Gamma law is
\[
 C_x(p)^4=\frac{(1+q_{x,p})^4}{8(1+q_{x,p}^4)},\qquad
 \frac12<C_x(p)<1.
 \tag{6.8}\label{eq:gamma-fourth}
\]
For every nonnegative prime energy \(a_p\),
\[
 \sum_p C_x(p)^2a_p<\infty
 \quad\Longleftrightarrow\quad
 \sum_p a_p<\infty.
 \tag{6.9}\label{eq:no-free-hilbertization}
\]
\item For every \(k,e\ge0\),
\[
 \lim_{x\downarrow0}
 \frac{(-1)^k\Delta^kF_{x,p}(e)}{x^{k+1}}
 =\Sigma_k(p):=\frac{q_p}{1+q_p}(\log p)^{k+1}.
 \tag{6.10}\label{eq:rees-symbol}
\]
If \(p_0,\ldots,p_{N-1}\) are distinct primes, then
\[
 \det\!\bigl[\Sigma_k(p_j)\bigr]_{0\le k,j<N}\ne0.
 \tag{6.11}\label{eq:rees-rank}
\]
\end{enumerate}
\end{theorem}

\begin{proof}
The identity \(b_{p,e}=\delta_pg_{p,e}\) is direct, and summing the geometric square gives
\[
 \sum_{e\ge0}g_{p,e}^2
 =\frac{q_p(1-q_p)^2}{1-q_p^2}
 =\frac{q_p(1-q_p)}{1+q_p}.
\]
The last expression is comparable to \(p^{-1}\), so its prime sum diverges.

Formula \eqref{eq:gamma-fourth} follows from
\(\cosh(2x\log p)=(q_{x,p}^{-2}+q_{x,p}^{2})/2\).
For \(0<q<1\),
\[
 8(1+q^4)-(1+q)^4=(1-q)^2(7q^2+10q+7)>0,
\]
while \(1+q^4<(1+q)^4\); positivity gives the stated bounds.  Hence
\(\tfrac14a_p<C_x(p)^2a_p<a_p\), which proves \eqref{eq:no-free-hilbertization}.

Iteration of the forward difference in \eqref{eq:gamma-ray} gives
\[
 \Delta^kF_{x,p}(e)
 =(-1)^k C_x(p)\frac{q_p}{1+q_p}
   (1-q_{x,p})^{k+1}q_{x,p}^e.
\]
As \(x\downarrow0\), \(C_x(p)\to1\),
\((1-p^{-x})/x\to\log p\), and \(p^{-xe}\to1\), proving
\eqref{eq:rees-symbol}.  Finally
\[
 \det\!\bigl[\Sigma_k(p_j)\bigr]
 =\left(\prod_{j=0}^{N-1}\frac{q_{p_j}}{1+q_{p_j}}\log p_j\right)
   \prod_{0\le i<j<N}(\log p_j-\log p_i),
\]
up to the harmless choice of Vandermonde orientation.  Every factor is nonzero.
\end{proof}

The theorem identifies the positive survivor of the finite-completion obstruction exposed by
the source disintegration.  Fisher whitening necessarily restores prime-harmonic size, and Gamma sewing
cannot hide that divergence because it is bounded above and below uniformly.  What Gamma retains
is not a new bounded scalar carrier but every finite prime packet across an unbounded tower of
Rees grades.  This is the source-side object that finite scalarizations were approximating.

\subsection{The odd orientation jet before squaring}

The strict contraction in \eqref{eq:gamma-fourth} has a canonical missing channel.  Put
\[
 D_x(p)=\sqrt{1-C_x(p)^2},\qquad
 J_x(p)=
 \begin{pmatrix}C_x(p)&-D_x(p)\\ D_x(p)&C_x(p)\end{pmatrix},
 \qquad
 \ell_x(p)=\frac{2}{\sqrt3}\frac{D_x(p)}x.
 \tag{6.J1}\label{eq:julia-lift}
\]
Let \(\mathcal H=\ell^2(\N_0;\R)\), \(k_p=(1,p^{-1},p^{-2},\ldots)\), and, for a finitely
supported prime packet \(d=(d_p)\), define
\[
 C(d)=\sum_p d_pk_p,\qquad
 L_m(d)=\sum_p d_p(\log p)p^{-m}k_p.
 \tag{6.J2}\label{eq:prime-clock}
\]

\begin{theorem}[Julia--Fisher orientation transgression]
\label{thm:odd-orientation}
The following hold.
\begin{enumerate}[label=\textup{(\roman*)}]
\item \(J_x(p)\) is the unique positively oriented plane isometry with first column
\((C_x(p),D_x(p))^\top\), and
\[
 \ell_x(p)\longrightarrow\log p\qquad(x\downarrow0).
 \tag{6.J3}\label{eq:julia-clock}
\]
\item For every \(m\ge0\), the graph of the prime source and its logarithmic clock is maximally
singular:
\[
 \overline{\{(C(d),L_m(d)):d\text{ finitely supported}\}}
 =\mathcal H\oplus\mathcal H.
 \tag{6.J4}\label{eq:dense-clock-graph}
\]
The Julia clocks obtained by replacing \(\log p\) with \(\ell_x(p)\) converge on every finite
packet to \(L_m\).  Along any \(x_j\downarrow0\), every family of bounded operators which
recovers all these packet clocks has norm tending to infinity.
\item On any complex Hilbert space, let
\[
 P_{\psi,\eta}=2^{-1/2}
 (|\eta\rangle\langle\psi|+|\psi\rangle\langle\eta|),\qquad
 \delta_e(X)=Xe-eX,
\]
and
\[
 \mathfrak W(e,T;\psi,\eta)
 =\operatorname{Tr}\!\left(\delta_e(P_{\psi,\eta})^*\delta_e(T)\right).
 \tag{6.J5}\label{eq:pointed-ward}
\]
Then
\[
 \mathfrak W(e,T;\psi,-\eta)=-\mathfrak W(e,T;\psi,\eta).
 \tag{6.J6}\label{eq:ward-odd}
\]
If \(e=e^*\), finite-rank cyclicity gives the exact transgression
\[
 \mathfrak W(e,T;\psi,\eta)
 =\operatorname{Tr}\!\left(P_{\psi,\eta}\,\delta_e(\delta_e(T))\right).
 \tag{6.J7}\label{eq:ward-curvature}
\]
\end{enumerate}
\end{theorem}

\begin{proof}
Equation \eqref{eq:gamma-fourth} gives \(C_x(p)^2+D_x(p)^2=1\), so
\(J_x(p)\) is orthogonal with determinant one.  In dimension two, the positive first column
fixes the second column and proves uniqueness.  Expanding \eqref{eq:gamma-fourth} at
\(p^{-x}=1\) gives
\[
 1-C_x(p)^2=\frac34x^2(\log p)^2+o(x^2),
\]
which proves \eqref{eq:julia-clock}.

For the graph assertion, write \(F_f(z)=\sum_{j\ge0}f_jz^j\).  A vector
\((f,g)\) orthogonal to the graph in \eqref{eq:dense-clock-graph} satisfies
\[
 F_f(p^{-1})+(\log p)p^{-m}F_g(p^{-1})=0
 \tag{6.J8}\label{eq:prime-germ}
\]
for every prime.  The recursion \(F_f(z)=f_0+zF_{Bf}(z)\), together with
\(p^{-1}\to0\), \(\log p\to\infty\), and \(p^{-1}\log p\to0\), removes the first \(m\)
coefficients and reduces \eqref{eq:prime-germ} to the case \(m=0\).  Dividing first by
\(\log p\), then by \(p^{-1}\), and iterating the same recursion forces every coefficient of
\(f\) and \(g\) to vanish.  The graph therefore has trivial orthogonal complement.
Packetwise Julia convergence follows from \eqref{eq:julia-clock}.  If bounded recoveries had
uniformly bounded norm, their limiting graph would lie in a closed cone
\(\{(u,v):\|v\|\le M\|u\|\}\), contradicting density because the graph closure contains every
\((0,v)\).

Finally \(P_{\psi,-\eta}=-P_{\psi,\eta}\); linearity of \(\delta_e\), adjoint, and trace proves
\eqref{eq:ward-odd}.  When \(e=e^*\), expanding the two commutators and using cyclicity for the
finite-rank factor \(P_{\psi,\eta}\) gives \eqref{eq:ward-curvature}.
\end{proof}

This theorem is the source-side orientation law.  The scalar contraction has only quadratic
contact with its lossless boundary; the added Julia channel restores the missing first jet.
The pointed Ward observable remains linear in that jet.  Squaring \(D_x\), taking a Hellinger
affinity, or passing to an unpointed Hilbert--Schmidt energy produces an even quotient.  Those
quotients can detect a defect but cannot decide its direction.

\begin{theorem}[Even-current/odd-orientation splitting]
\label{thm:parity-splitting}
Let \(X\) denote the completed Xi carrier as a meromorphic upper-end germ.  Then \(X^2\) is
hypertranscendental over \(\C(s)\), and
\[
 \C(s)\bigl\langle (\log X^2)'\bigr\rangle_{\rm diff}
 =
 \C(s)\bigl\langle (\log X)'\bigr\rangle_{\rm diff}.
 \tag{6.P1}\label{eq:square-current}
\]
Thus squaring retains the complete logarithmic-current differential field even though it
forgets the choice of sign lift.

More generally, write \(\Xi(z)=\xi(\tfrac12+iz)\), let \(E,E^\sharp\) be an ordered reflected
analytic channel pair, and put
\[
 \mathfrak a(z)=\frac{E(z)+E^\sharp(z)}{E(z)-E^\sharp(z)}.
 \tag{6.P2}\label{eq:odd-cayley}
\]
Channel exchange sends \(\mathfrak a\) to \(-\mathfrak a\).  For the source-selected completed
theta pair, on every regular point of the open theta strip,
\[
 \mathfrak a(z)=0\quad\Longleftrightarrow\quad\Xi(z)=0,
 \tag{6.P3}\label{eq:odd-cayley-divisor}
\]
and sharp reflection conjugates \(\mathfrak a\) to \(-\mathfrak a\).
Finally, if \(\varphi\) is differentiable and even near zero, then
\[
 \varphi'(0)=0.
 \tag{6.P4}\label{eq:even-first-jet}
\]
\end{theorem}

\begin{proof}
For a nonzero meromorphic germ,
\[
 (\log X^2)'=2(\log X)'.
\]
The nonzero constant \(2\) is invertible in the base field, proving
\eqref{eq:square-current}; a positive power of a hypertranscendental germ remains
hypertranscendental.  Swapping \(E\) and \(E^\sharp\) preserves the numerator of
\eqref{eq:odd-cayley} and reverses its denominator.  For the completed theta pair,
\[
 E(z)+E^\sharp(z)=4z^2\Xi(z),
\]
and at an upper-strip Xi zero the denominator is nonzero, proving
\eqref{eq:odd-cayley-divisor}; the sharp law follows from exchange of the two channels.
If \(\varphi(-z)=\varphi(z)\), differentiating at zero gives
\(-\varphi'(0)=\varphi'(0)\), hence \eqref{eq:even-first-jet}.
\end{proof}

This resolves the apparent contradiction between divisor recovery after squaring and loss of
orientation.  Evenization preserves support, multiplicity up to the known power factor, and the
logarithmic current; it removes only the odd lift which distinguishes the reflected sides.
Ramification and orientation therefore cannot be reconstructed from the same parity quotient.

\begin{corollary}[Finite-reader and finite-local-frame boundary]
\label{cor:finite-reader}
Every linear readout of the Xi logarithmic-derivative jet tower into a finite-dimensional complex
space annihilates a nonzero finite derivative packet.  Adding any finite Gibbs or finite spectral
clock does not remove this kernel.  Moreover, on no positive aperture line can the horizontal
derivative of \(|\xi|^2\) be represented by a finite sum of autocorrelation energies of
constant-coefficient differential filters.
\end{corollary}

\begin{proof}
The logarithmic derivative of Xi is constant-coefficient differentially transcendental because
it has an actual pole.  One more derivative vector than the target dimension gives a nonzero
kernel packet; finite Gibbs dressing has finite differential rank and does not change this
conclusion.  A finite constant-coefficient autocorrelation frame would have a polynomial
spectral density.  Analytic continuation would turn the horizontal derivative identity into an
entire real-shift coboundary for the Xi logarithmic flux.  Such a coboundary propagates every Xi
zero by the same positive real shift with multiplicity.  Iteration eventually leaves the
critical strip, contradicting the unconditional zero location \(0<\Rea\rho<1\).
\end{proof}

\subsection{Three finite-completion failures of one prime source}

The same critical activation is already infinite before reflection is considered.

\begin{theorem}[Critical activation singular-resolution law]
\label{thm:critical-resolution}
At \(q_p=p^{-1}\), the occupancy, fixed-depth tail, and Fisher-whitened Gamma masses satisfy
\[
 \sum_p\frac1{p+1}=\infty,\qquad
 \sum_p\frac1{p^{m-1}(p+1)}<\infty\quad(m\ge2),
\]
\[
 \sum_p\frac{p-1}{p(p+1)}=\infty.
\]
No nonzero polynomial vanishes on infinitely many nodes \(p^{-1}\).

At a divisor of multiplicity \(d\ge1\), let the local overlap and first score numerator be
\[
 B_d(t)=t^d,\qquad S_d(t)=dt^{d-1}.
\]
Then
\[
 \frac{|S_d(t)|^2}{|B_d(t)|^2}=\frac{d^2}{|t|^2}.
 \tag{6.12}\label{eq:gram-blowup}
\]
Consequently no uniformly bounded diagonal \(C(t)\) makes
\[
 \begin{pmatrix}|B_d(t)|^2&S_d(t)\\
                 S_d(t)&C(t)\end{pmatrix}
\]
positive semidefinite near \(t=0\).
\end{theorem}

\begin{proof}
The first and third summands are comparable to \(p^{-1}\); the middle summand is bounded by
\(p^{-m}\).  The summability statements follow from Euler's divergence of
\(\sum_pp^{-1}\) and convergence of \(\sum_nn^{-m}\) for \(m\ge2\).  Infinitely many
distinct reciprocal primes cannot be roots of a nonzero polynomial.  Formula
\eqref{eq:gram-blowup} is immediate.  Positivity of the displayed matrix would force
\(|S_d(t)|^2\le |B_d(t)|^2C(t)\), hence \(C(t)\ge d^2/|t|^2\), contradicting uniform
boundedness.
\end{proof}

The algebraic completion has an independent infinite frontier.  Let \(I\) be finite, let
\(A=\prod_{i\in I}A_i\) be a finite product of pointed sets, and let
\(y:\mathcal P\to A\) be tail-complete, where \(\mathcal P\) is the set of rational primes:
every state occurs beyond every finite set of primes.
Let \((b_z)_{z\in A}\) be the anchored tensor basis of \(\C^A\), and put
\[
 F_{r,z}(\tau)=\sum_{p\in\mathcal P}b_z(y(p))p^{i(r+1)\tau},
 \qquad r\ge0.
 \tag{6.13}\label{eq:ht-grid-carrier}
\]
These are upper-half-plane Dirichlet germs.  Same-prime coefficient fusion is
\[
 V_e(f)\star V_r(g)=V_{e+r+1}(fg).
 \tag{6.14}\label{eq:valuation-fusion}
\]

\begin{lemma}[Tail-echelon carrier independence]\label{lem:tail-echelon}
For every finite set of pairs \((r,z)\), the corresponding germs
\((F_{r,z})\) are jointly differentially independent over \(\C(\tau)\).
\end{lemma}

\begin{proof}
Let \(c_{r,z}\) be the arithmetic column supported on \(p^{r+1}\) with coefficient
\(b_z(y(p))\).  In the arithmetic-function algebra, prime and logarithmic derivations give
\[
 \partial_q L^kc_{r,z}(q^r)
 =(r+1)\bigl(-(r+1)\log q\bigr)^k b_z(y(q)),
\]
and the same expression vanishes below total prime-power degree \(r\).  This is the block form of
the leading-echelon calculation in Theorem~\ref{thm:omega}.  At fixed depth, the matrix
\((b_z(a))_{z,a\in A}\) is invertible because \((b_z)\) is a basis of \(\C^A\).
Tail completeness supplies every column \(a\) after deletion of any finite prime set.  Hence the
leading block has full rank on every cofinite tail.

If a nonzero differential polynomial relation existed, choose its least leading prime-power
degree and apply enough prime derivations to isolate that block.  Powers of \(\log q\) first
separate jet orders, and the invertible evaluation matrix then separates \(z\)-labels.  Its
leading coefficient must vanish, a contradiction; induction removes every block.  This proves
differential independence over \(\C\).

Adjoining \(\tau\) creates no relation.  After clearing a rational denominator, a proposed
relation over \(\C(\tau)\) is a finite polynomial in \(\tau\) with Dirichlet-germ coefficients.
Choose the least Dirichlet exponent in its nonzero coefficients.  Upper-end exponential
separation extracts a nonzero polynomial in \(\tau\), which cannot vanish identically.
\end{proof}

\begin{theorem}[Neutral-source hypertranscendental creation law]
\label{thm:ht-creation}
Over the embedded coordinate field \(E=\C(\tau)\), the complete family
\((F_{r,z})_{r\ge0,z\in A}\) is countably jointly differentially independent.  If
\[
 K_n=E\langle F_{r,z}:0\le r\le n,\ z\in A\rangle_{\rm diff},
\]
then
\[
 K_0\subsetneq K_1\subsetneq\cdots,
\]
and the complete layer \((F_{n+1,z})_{z\in A}\) is differentially independent over \(K_n\) of
exact rank \(|A|=\prod_i|A_i|\).  Nevertheless
\[
 V_n(b_z)\star V_0(1)=V_{n+1}(b_z),
 \qquad F_{n+1,z}\notin K_n.
 \tag{6.15}\label{eq:external-creation}
\]
Thus every derivation-closed field which contains depth zero and is closed under the exact
fusion \eqref{eq:valuation-fusion} contains the whole countably free family.

For pairwise coprime moduli \(q_i\), the actual prime profile
\(y(p)=(p\bmod q_i)_i\) in
\(\prod_i(\mathbb Z/q_i\mathbb Z)^\times\) is tail-complete, and each frontier has exact rank
\(\prod_i\varphi(q_i)\).  The all-identity sector is the prime-zeta spine
\[
 F_{r,\mathbf1}(\tau)=\sum_pp^{i(r+1)\tau}.
\]
\end{theorem}

\begin{proof}
Lemma~\ref{lem:tail-echelon}, applied to arbitrary finite packets, gives countable joint
differential independence.  Applying it to all layers through \(n+1\) shows that the last layer is
independent over the differential field generated by the complete prefix and gives its exact
cardinality.

The identity observable has empty support and satisfies \(b_z\cdot1=b_z\).
Equation \eqref{eq:valuation-fusion} therefore gives
\eqref{eq:external-creation}, while relative independence gives its nonmembership assertion.
Induction proves that any fusion-complete field containing depth zero contains every layer.
For the arithmetic specialization, the Chinese remainder theorem and Dirichlet's theorem on
primes in progressions give tail completeness; the rank is the product of the unit-group
cardinalities.
\end{proof}

Theorems~\ref{thm:critical-resolution} and \ref{thm:ht-creation} expose two different failures.
Finite differential closure cannot contain the source's own fusion orbit.  Bounded positive
completion cannot cross its divisor because differentiation lowers overlap grade by one.
The next theorem identifies the singular quotient forced by the second failure.

\subsection{Exact squarefree graph and the inner--outer weld}

Put
\[
 Q(s)=\frac{\zeta(s)}{\zeta(2s)},\qquad
 \mathcal A(s)=-\frac{\zeta'(s)}{\zeta(s)},\qquad
 C_{\rm sf}(s)=\mathcal A(s)-2\mathcal A(2s),
\]
and let
\[
 \gamma_\infty(s)=\frac{s(s-1)}2\pi^{-s/2}\Gamma(s/2).
\]
For \(s=\sigma+u\varepsilon\), define the moving overlap and its source score by
\[
 O_{\sigma,\varepsilon}(u)=\frac{Q(s)}{Q(\sigma)},\qquad
 S_{\sigma,\varepsilon}(u)
 =2\left[\partial_uO_{\sigma,\varepsilon}(u)
   +\varepsilon C_{\rm sf}(\sigma)O_{\sigma,\varepsilon}(u)\right].
\]
The singular squarefree graph is
\[
 \mathfrak G_{\rm sf}(\sigma,\varepsilon;u)
 =2C_{\rm sf}(\sigma)
  -\frac{S_{\sigma,\varepsilon}(u)}
         {\varepsilon O_{\sigma,\varepsilon}(u)}
  -2\frac{\gamma_\infty'(s)}{\gamma_\infty(s)}.
 \tag{6.16}\label{eq:sf-graph}
\]

\begin{theorem}[Exact graph/connection/nearby-profile weld]\label{thm:graph-weld}
If \(\Rea\sigma>1\), \(\Rea s>1\), and \(\varepsilon\ne0\), then
\[
 \mathfrak G_{\rm sf}(\sigma,\varepsilon;u)
 =2\left[C_{\rm sf}(s)
   -\frac{\gamma_\infty'(s)}{\gamma_\infty(s)}\right].
 \tag{6.17}\label{eq:sf-graph-current}
\]
After the forced strict-right correction,
\[
 \mathcal G_\xi(s):=
 \mathfrak G_{\rm sf}(\sigma,\varepsilon;u)+4\mathcal A(2s)
 =-2\frac{\xi'(s)}{\xi(s)}.
 \tag{6.18}\label{eq:xi-graph}
\]
Consequently
\[
 \frac{s-\tfrac12}{2}\mathcal G_\xi(s)=J_\xi(s).
 \tag{6.19}\label{eq:graph-euler}
\]
At every zero \(\rho\), the residue of the meromorphic right side of
\eqref{eq:graph-euler} equals the first outer coefficient of the balanced inner profile:
\[
 \Res_\rho J_\xi
 =-m_\rho(\rho-\tfrac12)
 =\lim_{w\to0,\ w\ne0}w^{-1}
 \left[
 \left(\frac{w^{-1}}
 {w^{-1}+\rho-\tfrac12}\right)^{m_\rho}-1
 \right].
 \tag{6.20}\label{eq:graph-inner-outer}
\]
\end{theorem}

\begin{proof}
Logarithmic differentiation of the overlap gives
\[
 \frac{\partial_uO_{\sigma,\varepsilon}(u)}
      {\varepsilon O_{\sigma,\varepsilon}(u)}
 =-C_{\rm sf}(s).
\]
Substitution in \eqref{eq:sf-graph} proves \eqref{eq:sf-graph-current}.
Since
\[
 C_{\rm sf}(s)+2\mathcal A(2s)=\mathcal A(s)
\quad\text{and}\quad
\xi=\gamma_\infty\zeta,
\]
equation \eqref{eq:xi-graph} follows, and
\eqref{eq:graph-euler} is its balanced Euler contraction.  Equation
\eqref{eq:graph-inner-outer} is Theorem~\ref{thm:inner-outer}.
\end{proof}

The correction \(4\mathcal A(2s)\) is holomorphic on \(\Rea s>\tfrac12\); it changes the
squarefree current into the completed Xi current without adding a right-half divisor.  More
importantly, division by \(O_{\sigma,\varepsilon}\) is forced by
\eqref{eq:gram-blowup}.  The graph is a source-defined singular correspondence, not a bounded
Fisher tangent disguised by normalization.

For a prime \(p\) and \(c<\tfrac12\), let
\[
 \mu_{p,c}(k)=(1-r_{p,c})r_{p,c}^{\,k},
 \qquad r_{p,c}=p^{-1+2c},\qquad k\ge0.
 \tag{6.21}\label{eq:geom-law}
\]
This is exactly \eqref{eq:conditional-depth} with
\(\sigma=1-2c\), after shifting positive valuation \(k=e+1\) to \(e\ge0\).
Critical reflection \(c\mapsto-c\) therefore compares the two conditional valuation fibres at
the exponents \(1-2c\) and \(1+2c\).

Compare \(\mu_{p,c}\) with the critical law \(\mu_{p,0}\).  Its squared Hellinger affinity is
\[
 h_p(c)^2
 =\frac{(1-p^{-1})(1-p^{-1+2c})}
        {(1-p^{-1+c})^2}.
 \tag{6.22}\label{eq:local-hellinger}
\]
For a prime cutoff \(N\), put
\[
 H_N(c)=\prod_{p<N}h_p(c),\qquad
 a_N=\prod_{p<N}(1-p^{-1}).
 \tag{6.23}\label{eq:HN-aN}
\]
The scale \(a_N\to0\) is the Wold vacuum mass.

\begin{definition}[All-orders flatness]
A nonnegative germ \(x_N\) is \emph{flat beyond every Rees order} relative to \(a_N\) if
\[
 x_N=o(a_N^K)\qquad\text{for every }K\in\N.
 \tag{6.24}\label{eq:flat}
\]
\end{definition}

\begin{theorem}[Hellinger--Rees trichotomy]\label{thm:trichotomy}
For \(c<\tfrac12\):
\[
 \begin{array}{rcll}
 c<0&:&H_N(c)=\Theta(a_N^{1/2}),&\text{half-normal face},\\[2mm]
 c=0&:&H_N(c)=1,&\text{critical face},\\[2mm]
 0<c<\tfrac12&:&H_N(c)=o(a_N^K)\quad\forall K,&\text{all-orders-flat face}.
 \end{array}
 \tag{6.25}\label{eq:trichotomy}
\]
\end{theorem}

\begin{proof}
The middle line is immediate.  Suppose \(c<0\).  Since
\(p^{-1+2c}\le p^{-1+c}\le p^{-1}\), the local factor
\eqref{eq:local-hellinger} is bounded below by \(1-p^{-1}\).  Conversely,
\[
 h_p(c)^2\le(1-p^{-1})\exp(C_c p^{-1+c})
\]
with a constant \(C_c\) independent of \(p\).  The correction product converges because
\(\sum_pp^{-1+c}<\infty\).  Hence \(H_N(c)^2=\Theta(a_N)\).

Now suppose \(c>0\).  For every fixed \(K\), the ratio
\(p^{-1+2c}/p^{-1}=p^{2c}\) tends to infinity.  Direct subtraction in
\eqref{eq:local-hellinger} gives
\[
 1-h_p(c)^2
 =\frac{p^{-1}(p^c-1)^2}{(1-p^{-1+c})^2}.
\]
For all sufficiently large \(p\), this is at least
\(\tfrac14p^{-1+2c}\), hence at least \((2K+2)p^{-1}\).
Bernoulli's inequality therefore yields
\[
 h_p(c)^2\le(1-p^{-1})^{2K+2}.
\]
The finite exceptional product contributes only a constant, so
\[
 H_N(c)^2=O(a_N^{2K+2})=o(a_N^{2K}).
\]
Taking square roots gives \(H_N(c)=o(a_N^K)\).
\end{proof}

\begin{remark}
Both off-critical faces have \(H_N(c)\to0\), so the infinite products are mutually singular by
Kakutani's theorem.  The ordinary singular/nonsingular dichotomy identifies the two faces.
Equation \eqref{eq:trichotomy} resolves the orientation discarded by that quotient.
\end{remark}

\section{Universal asymptotic superselection}\label{sec:superselection}

Let \(E,F\) be real normed spaces and \(a_N>0\) with \(a_N\to0\).

\begin{theorem}[Essential sewing cost]\label{thm:essential}
Suppose
\[
 \|x_N\|=o(a_N^m)\quad\text{for every }m,\qquad
 \|y_N\|=\Theta(a_N^d)
 \tag{7.1}\label{eq:types}
\]
for some \(d\in\N\), with \(x_N\ne0\) eventually.  If bounded linear operators
\(R_N:E\to F\) satisfy
\[
 R_Nx_N=y_N,
 \tag{7.2}\label{eq:sewing}
\]
then
\[
 a_N^{-K}=o(\|R_N\|)
 \qquad\text{for every }K\in\N.
 \tag{7.3}\label{eq:essential-cost}
\]
\end{theorem}

\begin{proof}
Fix \(K\).  Flatness with order \(K+d\) and finite grade \(d\) give
\[
 \frac{\|x_N\|}{a_N^{K+d}}\longrightarrow0,\qquad
 \|y_N\|\asymp a_N^d.
\]
Therefore
\[
 a_N^{-K}=o\!\left(\frac{\|y_N\|}{\|x_N\|}\right).
\]
The operator norm inequality and \eqref{eq:sewing} imply
\[
 \frac{\|y_N\|}{\|x_N\|}\le\|R_N\|,
\]
which proves \eqref{eq:essential-cost}.
\end{proof}

\begin{corollary}[Flat-radical invariance]\label{cor:flat-invariant}
If \(\|R_N\|=O(a_N^{-K})\) for one finite \(K\), then \(R_N\) maps every all-orders-flat germ
to an all-orders-flat germ.  The conclusion holds in arbitrary dimension and for finite matrices
whose entries share one finite Rees bound.
\end{corollary}

\begin{proof}
For every \(m\), apply the operator bound to flatness at order \(m+K\).
\end{proof}

\begin{corollary}[Xi-charge superselection]\label{cor:xi-superselection}
Let \(\rho\) be a right-half zero of \(\xi\), and put
\(c_\rho=\Rea\rho-\tfrac12\).  No finite-Rees-moderate operator family sews
\[
 H_N(c_\rho)\longmapsto H_N(-c_\rho).
 \]
Every exact sewing family has norm satisfying \eqref{eq:essential-cost}.
\end{corollary}

\begin{proof}
The critical strip gives \(0<c_\rho<\tfrac12\).
Theorem~\ref{thm:trichotomy} makes the source germ all-orders flat and the target half-normal,
so Corollary~\ref{cor:flat-invariant} already rules out finite-Rees sewing.  For the quantitative
claim, apply Theorem~\ref{thm:essential} to the squared affinities:
\[
 H_N(c_\rho)^2=o(a_N^m)\quad(\forall m),\qquad
 H_N(-c_\rho)^2=\Theta(a_N).
\]
It follows that
\[
 a_N^{-2K}
 =o\!\left(\frac{H_N(-c_\rho)^2}{H_N(c_\rho)^2}\right).
\]
All quantities are positive; taking square roots gives
\[
 a_N^{-K}
 =o\!\left(\frac{H_N(-c_\rho)}{H_N(c_\rho)}\right).
\]
The operator-norm inequality finishes the proof.
\end{proof}

\begin{remark}
This is stronger than saying that the two faces have different finite grades.  One face lies in
the intersection of every positive Rees grade.  That intersection is an invariant radical.
Increasing matrix size, Hilbert-space dimension, or finite jet order does not cross it.  Exact
sewing requires a singularity beyond every finite Rees order.
\end{remark}

\section{Why local scalar observers miss the same charge}\label{sec:optimality}

Functional-equation-preserving perturbative instability is known
\cite{GauthierXarles2009}.  The boundary character is optimal in the following stronger,
quantified compact-open sense.

\begin{theorem}[Escaping-quartet separation]\label{thm:escape}
Fix \(0<c<\tfrac12\).  For \(T>0\), set
\[
 G_{c,T}(s)=
 \left(1+\frac{((s-\tfrac12)-c)^2}{T^2}\right)
 \left(1+\frac{((s-\tfrac12)+c)^2}{T^2}\right),
 \qquad F_{c,T}=G_{c,T}\xi.
 \tag{8.1}\label{eq:quartet-gauge}
\]
Then:
\begin{enumerate}[label=\textup{(\roman*)}]
\item \(F_{c,T}\) is entire, obeys both completed symmetries, is hypertranscendental, and has the
added quartet \(\tfrac12\pm c\pm iT\);
\item as \(T\to\infty\), \(F_{c,T}\) and every fixed derivative converge locally uniformly to
\(\xi\);
\item the nonnegative Ward density
\[
 \mathcal W_{c,T}(t)=
 2\left[
 \frac{2(c+\tfrac12)}{(c+\tfrac12)^2+(t+T)^2}
 \frac{2(c+\tfrac12)}{(c+\tfrac12)^2+(t-T)^2}
 \right]
 \tag{8.2}\label{eq:ward-density}
\]
converges uniformly to zero on every bounded interval, but
\[
 \int_{\R}\mathcal W_{c,T}(t)\,\dd t=8\pi;
 \tag{8.3}\label{eq:ward-mass}
\]
\item the reciprocal-square energy
\[
 \mathcal E_{c,T}
 =\frac{\sin^2(2\pi c)}{c^2+T^2}
 \tag{8.4}\label{eq:trace-energy}
\]
tends to zero, whereas the centered character of
\(\rho_T=\tfrac12+c+iT\) satisfies, for every real \(t\),
\[
 \left|e^{t(\rho_T-1/2)}\right|=e^{tc},
 \qquad
 \frac{|\arg e^{-2\pi i(\rho_T-1/2)}|}{2\pi}=c.
 \tag{8.5}\label{eq:escape-character}
\]
\end{enumerate}
\end{theorem}

\begin{proof}
The polynomial \(G_{c,T}\) is invariant under \(s\mapsto1-s\) and conjugation, and its four zeros
are the displayed quartet.  It is nonzero as a polynomial.  To recall why \(\xi\) is
hypertranscendental, put \(L=\xi'/\xi\).  On \(\Rea s>1\), the Gamma recurrence and Euler series
give
\[
 L(s+2)-L(s)
 -\left(\frac1{s+2}+\frac1{s+1}-\frac1{s-1}\right)
 =\sum_{n\ge1}\frac{\Lambda(n)(1-n^{-2})}{n^s}.
\]
The coefficient at every prime \(p\) is \(\log p(1-p^{-2})\ne0\), while the support is contained
in the prime powers.  Ostrowski's Dirichlet-series theorem makes the right side
hypertranscendental \cite{Ostrowski1920,Nagoshi2016}.  Differential algebraicity of \(\xi\)
would pass to \(L\), its translate, and their rationally corrected difference, a contradiction.
If \(F_{c,T}\) were differentially algebraic, division by
\(G_{c,T}\in\C(s)^\times\) would now make \(\xi\) differentially algebraic.  On each compact
set, \(G_{c,T}\to1\) together with every fixed derivative; Leibniz's rule proves the asserted jet
convergence.

Put \(a=c+\tfrac12>0\).  Each summand \(2a/(a^2+(t\pm T)^2)\) is nonnegative and has integral
\(2\pi\).  This proves \eqref{eq:ward-mass}.  If \(|t|\le R<T\), then
\[
 0\le\mathcal W_{c,T}(t)\le\frac8{(T-R)^2},
\]
which proves uniform compact-window escape.  Equation \eqref{eq:trace-energy} tends to zero
directly.  Finally, taking the norm and the principal argument of the exponential gives
\eqref{eq:escape-character}; the condition \(0<c<1/2\) prevents phase aliasing.
\end{proof}

Thus compact holomorphic jets, bounded Ward windows, and reciprocal energies can simultaneously
lose a fixed off-critical displacement.  The asymptotic character retains it exactly.  This does
not show that any named spectral program must fail.  It proves an information boundary: any
criterion continuous only in those local reporters cannot distinguish the escaping quartet.  It
must include a noncompact character coordinate or an equivalent singular boundary datum.

\section{The ramified-character boundary law}\label{sec:equalizer}

The preceding results still leave several apparently external phenomena: the theta dense
radical, analytic heat-node nonseparability, Wold--Poisson ancestry, Sonin/Calkin compactness,
four-unit Ward positivity, dyadic root descent, and local deck charge.  They become one phase
diagram only after continuous reflection character and discrete ramification are separated.

\subsection{The vertical axis: a divisor detector is necessarily singular}

We first record the topological mechanism without zeta functions.

\begin{proposition}[Cyclic-null collapse]\label{prop:cyclic-null}
Let \(\mathcal D\) be a complex vector space, let \(j:\mathcal D\to\mathcal H\) be an injective
linear map into a Hilbert space, and let \(Q\) be a Hermitian sesquilinear form on
\(\mathcal D\).  Suppose a family \((v_i)_{i\in I}\) lies in the left radical of \(Q\) and
\[
 \overline{\operatorname{span}\{j(v_i):i\in I\}}=\mathcal H.
 \tag{9.1}\label{eq:cyclic-null}
\]
Then for every \(u\in\mathcal D\) there is a sequence \(z_n\in\mathcal D\) such that
\[
 j(z_n)\longrightarrow0,\qquad
 Q(z_m-z_n,z_m-z_n)=0,\qquad
 Q(z_n,z_n)=Q(u,u).
 \tag{9.2}\label{eq:constant-energy-null}
\]
If \(Q\ne0\), its radical image is dense, proper, and nonclosed in \(\mathcal H\), and \(Q\)
is not sequentially closable through \(j\).
\end{proposition}

\begin{proof}
Choose \(r_n\) in the algebraic span of the \(v_i\) with \(j(r_n)\to j(u)\), and put
\(z_n=u-r_n\).  Hermitian symmetry makes the left radical two-sided.  Expansion gives
\eqref{eq:constant-energy-null}.  Cyclicity makes the radical image dense; injectivity and
\(Q\ne0\) make it proper, hence nonclosed.  Choosing \(u\) with \(Q(u,u)\ne0\) makes
\eqref{eq:constant-energy-null} a form-Cauchy sequence converging to zero in \(\mathcal H\) with
nonvanishing energy, which is exactly failure of sequential closability.
\end{proof}

Put \(\Xi(t)=\xi(\tfrac12+it)\), and let \(k\in\mathcal S(\R)\) be the completed theta source
normalized by \(\widehat k=\Xi\).  Hedenmalm's theta pair pulls back to momentum symmetry and
theta nullity for a source-core form \cite{Hedenmalm2026}.

\begin{theorem}[Theta-divisor rigging classification]\label{thm:theta-rigging}
Let \(Q:\mathcal S(\R)\times\mathcal S(\R)\to\C\) be jointly Schwartz-continuous,
sesquilinear, nonnegative, and momentum-symmetric:
\[
 Q(-i\partial_xf,g)=Q(f,-i\partial_xg).
 \tag{9.3}\label{eq:momentum-symmetry}
\]
If \(Q(k,k)=0\), there is a unique positive Borel measure \(\mu\) of temperate growth such that
\[
 Q(f,g)=\int_{\R}\widehat f(t)\overline{\widehat g(t)}\,\dd\mu(t),
 \qquad
 \supp\mu\subseteq\{t\in\R:\Xi(t)=0\}.
 \tag{9.4}\label{eq:theta-divisor-measure}
\]
The measure is atomic.  Conversely every positive temperate measure with this support defines
such a form.  If \(Q\ne0\), its ordinary-\(L^2\) radical image is dense and nonclosed and \(Q\)
is not sequentially closable through
\(\mathcal S(\R)\hookrightarrow L^2(\R)\).  In particular, within this source-core class,
\[
 Q\text{ is closable through }L^2(\R)
 \quad\Longleftrightarrow\quad Q=0.
 \tag{9.5}\label{eq:theta-closable-zero}
\]
\end{theorem}

\begin{proof}
Equation \eqref{eq:momentum-symmetry} makes \(Q\) invariant under simultaneous translation.
The Bochner--Schwartz theorem gives \eqref{eq:theta-divisor-measure}
\cite{Schwartz1966}.  Nonnegativity supplies Cauchy--Schwarz, so \(Q(k,k)=0\) puts \(k\) in both
radicals.  Hence
\[
 0=Q(k,k)=\int_{\R}|\Xi(t)|^2\,\dd\mu(t),
\]
which forces the stated support.  Real zeros of the nonzero entire function \(\Xi\) are
discrete, so \(\mu\) is atomic; the converse is immediate.

Every real translate of \(k\) lies in the radical.  Since \(\Xi\) is nonzero almost everywhere
on \(\R\), Wiener's theorem makes those translates cyclic in ordinary \(L^2\)
\cite{Wiener1932}.  Proposition~\ref{prop:cyclic-null} gives the dense radical,
nonclosedness, and nonclosability.  The zero form is closable, proving
\eqref{eq:theta-closable-zero}.
\end{proof}

This is the vertical obstruction.  It persists on the real Xi divisor, where horizontal
character rate is already zero.  It therefore cannot be explained as an off-line orientation
defect.

\begin{theorem}[Analytic heat-node nonseparability]\label{thm:heat-node}
Put \(x_m=e^{-\pi m^2}\) for \(m\ge1\).  If an entire function \(F\) vanishes at every
\(x_m\) except possibly \(x_n\), then \(F(x_n)=0\) as well.  In particular no entire symbol
realizes the Kronecker data
\[
 F(x_n)=1,\qquad F(x_m)=0\quad(m\ne n).
 \tag{9.H}\label{eq:heat-node-delta}
\]
The same conclusion holds for a function holomorphic on any connected domain containing zero
and every \(x_m\).
\end{theorem}

\begin{proof}
The positive nodes \(x_m\) converge to zero.  If the zeros \(F(x_m)=0\), \(m\ne n\), accumulate
at zero inside a connected holomorphy domain, the identity theorem makes \(F\) identically zero
there.  Hence \(F(x_n)=0\), contradicting \eqref{eq:heat-node-delta}.
\end{proof}

Theorems~\ref{thm:theta-rigging} and \ref{thm:heat-node} are complementary.  Ordinary \(L^2\)
closure collapses a nonzero divisor form because the theta source is cyclic; analytic
compactification cannot isolate one source node because the nodes accumulate at the boundary.
The surviving detector must therefore be singular in both Hilbert and analytic categories.

\subsection{The complete divisor output: continuous character and ramification}

Define the critical reflection
\[
 \rho^\dagger=1-\overline\rho
\]
and attach to a divisor point of multiplicity \(m_\rho\)
\[
 \chi_\rho(z)=\exp\!\bigl(z(\rho-\tfrac12)\bigr),\qquad
 \omega_{\rho,k}=
 \exp\!\left(\frac{2\pi i\,m_\rho}{2^k}\right).
 \tag{9.R1}\label{eq:ramified-character}
\]

\begin{theorem}[Ramified divisor-character law]\label{thm:ramified-character}
For every divisor point \(\rho\), every \(z,w\in\C\), every \(t\in\R\), and every \(k\in\N\),
\[
 \chi_\rho(z+w)=\chi_\rho(z)\chi_\rho(w),\qquad
 |\chi_\rho(t)|=e^{t(\Rea\rho-1/2)},
 \tag{9.R2}\label{eq:character-law}
\]
\[
 \chi_{\rho^\dagger}(t)=\overline{\chi_\rho(t)}^{-1},\qquad
 t\ne0\Longrightarrow
 \bigl(|\chi_\rho(t)|=1\Longleftrightarrow\Rea\rho=\tfrac12\bigr),
 \tag{9.R3}\label{eq:character-reflection}
\]
and
\[
 \omega_{\rho,k}^{\,2^k}=1,\qquad
 \omega_{\rho,k}=1\Longleftrightarrow2^k\mid m_\rho.
 \tag{9.R4}\label{eq:deck-divisibility}
\]
The continuous and ramified coordinates meet in the exact inner--outer charge
\[
 -m_\rho(\rho-\tfrac12).
 \tag{9.R5}\label{eq:ramified-infinitesimal}
\]
It is simultaneously the residue of the outer source-generated connection and the first
coefficient of the inner two-endpoint nearby profile.
\end{theorem}

\begin{proof}
Equations \eqref{eq:character-law}--\eqref{eq:character-reflection} follow from the exponential
addition law and
\[
 \rho^\dagger-\tfrac12=-\overline{\rho-\tfrac12}.
\]
For nonzero \(t\), \(e^{t(\Rea\rho-1/2)}=1\) exactly when
\(\Rea\rho=\tfrac12\).  Equation \eqref{eq:deck-divisibility} is the elementary root-of-unity
criterion for \(e^{2\pi i m_\rho/2^k}\).  Equation
\eqref{eq:ramified-infinitesimal} is Theorem~\ref{thm:graph-weld}.
\end{proof}

The two coordinates retain independent divisor data.  The full function
\(\chi_\rho(z)=e^{z(\rho-1/2)}\) carries the complex exponent \(\rho-\tfrac12\), while
\(m_\rho\) retains the divisor copy number.  The real norm of
\(\chi_\rho\) sees displacement but not multiplicity; the deck tower sees divisibility of
multiplicity but not displacement.  The residue \eqref{eq:ramified-infinitesimal} is their first
coupling, not a replacement for either coordinate.

\subsection{Quadratic compression detects the equalizer but erases its side}

The continuous character admits a canonical finite positive compression.  For a Xi divisor
copy \(\rho\), put
\[
 q_\chi(\rho)
 =\frac{\Im(\chi_\rho(-2\pi i))^2}{|\chi_\rho(-2\pi i)|^2}
 =\sin^2(2\pi\Rea\rho),
 \tag{9.Q1}\label{eq:character-defect}
\]
and define
\[
 \mathcal E_\chi
 =\sum_{\xi(\rho)=0}
 \frac{m_\rho}{|\rho-\tfrac12|^2}\,q_\chi(\rho),
 \tag{9.Q2}\label{eq:character-energy}
\]
where the sum is over divisor points and \(m_\rho\) is counted once in the displayed weight.

\begin{theorem}[Faithful quadratic character compression]
\label{thm:character-energy}
The series \eqref{eq:character-energy} converges absolutely and
\[
 \mathcal E_\chi\ge0.
 \tag{9.Q3}\label{eq:energy-positive}
\]
Moreover,
\[
 \mathcal E_\chi=0
 \quad\Longleftrightarrow\quad
 |\chi_\rho(1)|=1\ \text{for every Xi divisor point }\rho
 \quad\Longleftrightarrow\quad
 \Rea\rho=\tfrac12\ \text{for every such }\rho.
 \tag{9.Q4}\label{eq:energy-equalizer}
\]
If the last condition fails, a single off-line divisor point contributes a strictly positive
explicit lower bound to \(\mathcal E_\chi\).  Nevertheless
\[
 q_\chi(1-\overline\rho)=q_\chi(\rho),
 \tag{9.Q5}\label{eq:energy-even}
\]
so the compression is faithful to the zero locus of the character rate but not to its sign.
\end{theorem}

\begin{proof}
The centered reciprocal-square divisor moment
\(\sum_\rho m_\rho|\rho-\tfrac12|^{-2}\) is finite, while
\(0\le q_\chi(\rho)\le1\); comparison proves convergence and
\eqref{eq:energy-positive}.  Every weight is strictly positive.  Hence the sum vanishes exactly
when every \(q_\chi(\rho)\) vanishes.  Since \(0<\Rea\rho<1\),
\(\sin(2\pi\Rea\rho)=0\) exactly at \(\Rea\rho=\tfrac12\).  The character norm identity
\(|\chi_\rho(1)|=e^{\Rea\rho-1/2}\) gives the middle equivalence.  Retaining any one positive
summand gives the witness lower bound.  Finally
\(\sin(2\pi(1-\Rea\rho))=-\sin(2\pi\Rea\rho)\), proving
\eqref{eq:energy-even}.
\end{proof}

This is the exact status of scalar positivity in the architecture.  The energy is a complete,
trace-class detector of failure of the horizontal equalizer; it is not a mechanism that forces
the energy to vanish.  Equation \eqref{eq:energy-even} explains why a proliferation of positive
RH criteria can remain orientation-blind.  The unsquared source tangent in
Theorem~\ref{thm:odd-orientation} and the full character \(\chi_\rho\), not their quadratic
compression, carry the missing one-sided datum.

\begin{theorem}[Conservative character descent]\label{thm:conservative-descent}
Let \(E,F\) be real normed spaces, let \(U:E\to E\), \(V:F\to F\), and \(T:E\to F\) be
linear, and suppose
\[
 T U=V T,\qquad Ux=e^\nu x,\qquad Tx\ne0.
\]
If \(V\) is norm-conservative, meaning \(\|Vy\|=\|y\|\) for every \(y\in F\), then
\(\nu=0\).
\end{theorem}

\begin{proof}
Intertwining gives \(V(Tx)=e^\nu Tx\).  Taking norms and using \(Tx\ne0\) gives
\[
 \|Tx\|=\|V(Tx)\|=e^\nu\|Tx\|,
\]
so \(e^\nu=1\), hence \(\nu=0\).
\end{proof}

This theorem is not specific to zeta or Hilbert space.  It identifies the common logic of every
faithful conservative consumer: a nonunitary real character cannot survive it.  Off the
equalizer, either \(T\) kills the character vector or the target action fails conservativity.
The next results quantify the second alternative in four different categories.

\subsection{The continuous axis: one rate, four regular consumers}

For \(0<\Rea\rho<1\), put
\[
 \nu(\rho)=2\Rea\rho-1.
 \tag{9.6}\label{eq:boundary-rate}
\]
Theorem~\ref{thm:duality-fracture} says that \(\nu(\rho)=\Delta(\rho)\).  The same number is the
transport rate of the independent Wold-depth and Poisson-time coordinates.

Let \(M_n=\Rea H_n(1)>0\), where \(H_n\) is the shell multiplier from
\eqref{eq:shell-factor}, and define
\[
 \mathcal E_n(\rho;\tau)
 =\frac{|H_n(\rho)|^2}{M_n}e^{\nu(\rho)\tau}.
 \tag{9.7}\label{eq:two-clock-energy}
\]

\begin{theorem}[Wold--Poisson characteristic and quadrant law]\label{thm:wold-poisson}
For every \(\Rea\rho>0\) and \(u\in\R\),
\[
 \mathcal E_n(\rho;n\log2+u)
 \longrightarrow
 2|1-2^{-\rho}|^2e^{\nu(\rho)u}.
 \tag{9.8}\label{eq:two-clock-characteristic}
\]
For an arbitrary real rate \(\nu\), the conjugate two-clock character
\[
 C_\nu(\tau,d)=e^{\nu(\tau-d)}
 \tag{9.9}\label{eq:two-clock-character}
\]
is uniformly bounded on the independent quadrant \(\tau,d\ge0\) if and only if \(\nu=0\).
\end{theorem}

\begin{proof}
The least-atom and source-mass asymptotics are
\[
 2^{n\rho}H_n(\rho)\to1-2^{-\rho},
 \qquad
 2^nM_n\to\frac12.
 \]
Since \(e^{\nu(\rho)n\log2}=2^{n(2\Rea\rho-1)}\), substitution gives
\eqref{eq:two-clock-characteristic}.  On a characteristic line
\(\tau=d+u\), equation \eqref{eq:two-clock-character} equals \(e^{\nu u}\) for every \(\nu\).
If \(\nu>0\), it diverges along \((\tau,d)=(x,0)\); if \(\nu<0\), it diverges along
\((0,x)\).  For \(\nu=0\) it is identically one.
\end{proof}

Thus a comoving characteristic retains a finite nonzero profile even off the critical line.
The failed ancestry-purity argument had collapsed the two clocks before testing their independent
quadrant.  The missing information was not ancestry itself but the transverse character rate.

The same rate controls the Calkin quotient.  Burnol's Sonine/co-Poisson program attaches
Hilbert-space systems to the Riemann zeros and their multiplicities
\cite{Burnol2004}.  At positive Sonin aperture \(a\), let
\(\mathcal E_{p,q}^{(a)}\) be the ordered two-prime Euler Gram corner.  Its exact essential norm is
\[
 \|\mathcal E_{p,q}^{(a)}\|_{\rm ess}
 =(\sqrt p\sqrt q)^{-1}.
 \tag{9.10}\label{eq:sonin-corner-norm}
\]
Weight it by the prime seam from Theorem~\ref{thm:duality-fracture}:
\[
 \mathcal K_{p,q}^{(a)}(\rho)
 =(\kappa_p(\rho)-\kappa_q(\rho))\mathcal E_{p,q}^{(a)}.
 \tag{9.11}\label{eq:seam-calkin-corner}
\]

\begin{theorem}[Sonin--Calkin equalizer]\label{thm:calkin-equalizer}
For every pair of primes \(p,q\), every \(a>0\), and every
\(0<\Rea\rho<1\),
\[
 \|\mathcal K_{p,q}^{(a)}(\rho)\|_{\rm ess}
 =|\kappa_p(\rho)-\kappa_q(\rho)|(\sqrt p\sqrt q)^{-1}.
 \tag{9.12}\label{eq:seam-essential-norm}
\]
Consequently
\[
 \bigl[\mathcal K_{p,q}^{(a)}(\rho)\text{ is compact for all }p,q\bigr]
 \quad\Longleftrightarrow\quad
 \nu(\rho)=0.
 \tag{9.13}\label{eq:calkin-equalizer}
\]
\end{theorem}

\begin{proof}
Essential norm is homogeneous, so \eqref{eq:sonin-corner-norm} gives
\eqref{eq:seam-essential-norm}.  All corners are compact exactly when all prime seams are equal.
Theorem~\ref{thm:duality-fracture} proves that prime coherence is equivalent to
\(\Rea\rho=\tfrac12\), which is \(\nu(\rho)=0\).
\end{proof}

Finally let \(\mathcal D_R^+\) contain one upper Schwarz representative
\(\lambda=(\rho-\tfrac12)\) from every strict-right Xi zero pair and put, for \(\Rea w>0\),
\[
 \mathcal R_\Xi(w)
 =2\sum_{\lambda\in\mathcal D_R^+}
 \left(\frac1{w+\lambda}+\frac1{w+\overline\lambda}\right).
 \tag{9.14}\label{eq:ward-response}
\]
The Schwarz-paired sum is the continued anti-causal Ward response.  For a finite packet
\(w_i=\tfrac12+it_i\), let
\[
 G_{ij}=\frac1{w_i+\overline{w_j}},\qquad
 \mathcal W(t_1,\ldots,t_n)=
 \begin{pmatrix}
 G&(\mathcal R_\Xi(w_i))_i\\
 (\mathcal R_\Xi(w_i))_i^*&4
 \end{pmatrix}.
 \tag{9.15}\label{eq:ward-boundary-gram}
\]

\begin{theorem}[Ward equalizer and one-node witness]\label{thm:ward-equalizer}
The following global conditions are equivalent:
\[
 \begin{aligned}
 &\Delta(\rho)=0\quad\text{for every nontrivial Xi zero};\\
 &\mathcal W(t_1,\ldots,t_n)\succeq0
   \quad\text{for every finite real packet};\\
 &|\mathcal R_\Xi(\tfrac12+it)|^2\le4
   \quad\text{for every }t\in\R.
 \end{aligned}
 \tag{9.16}\label{eq:ward-equalizer}
\]
More locally, if \(\lambda=c+i\gamma\in\mathcal D_R^+\), then
\[
 \Rea\mathcal R_\Xi(\tfrac12-i\gamma)>2,\qquad
 |\mathcal R_\Xi(\tfrac12-i\gamma)|^2>4,
 \tag{9.17}\label{eq:ward-local-failure}
\]
so the corresponding one-node matrix in \eqref{eq:ward-boundary-gram} is not positive
semidefinite.
\end{theorem}

\begin{proof}
If every fracture vanishes, there are no strict-right zeros, the sum
\eqref{eq:ward-response} is empty, and the matrices are positive.  Conversely, at
\(w=\tfrac12-i\gamma\), the selected pair contributes a real part strictly larger than \(2\):
the term with \(w+\lambda=c+\tfrac12\) already contributes
\(2/(c+\tfrac12)>2\), while its Schwarz mate and every other pair have positive real part.
This proves \eqref{eq:ward-local-failure}.  At this node \(G_{11}=1\), and the Schur complement
of \eqref{eq:ward-boundary-gram} is \(4-|\mathcal R_\Xi(w)|^2<0\).  Hence any strict-right zero
violates both global conditions.  Functional-equation and Schwarz symmetry turn every off-line
zero into such an upper strict-right representative.
\end{proof}

\subsection{The ramification axis: dyadic roots and local loop charge}

For \(k\ge0\), let \(\mathsf D_k\) be the assertion that \(2^k\) divides the multiplicity of every
right-half Xi zero.

\begin{theorem}[Dyadic descent and sharp ramification level]\label{thm:dyadic-descent}
For every \(k\), the following are equivalent:
\begin{enumerate}[label=\textup{(\alph*)}]
\item \(\mathsf D_k\);
\item \(\xi\) has an analytic \(2^k\)-th root on \(\{\Rea s>\tfrac12\}\);
\item the functional-equation descent of \(\xi\) has an analytic \(2^k\)-th root on its slit
plane;
\item every sufficiently small circle about every right-half zero has trivial
\(\mu_{2^k}\)-deck character;
\item every finite height window admits compatible left/right root charts with trivial deck
transport;
\item a Schwarz-equivariant \(2^k\)-th root exists off the critical line.
\end{enumerate}
At a fixed right-half zero, every sufficiently small descended loop has deck character
\(\omega_{\rho,k}\).  In particular its rung-one character is persistently nontrivial exactly
when \(m_\rho\) is odd.

If RH is false and
\[
 H=\min_{\substack{\xi(\rho)=0\\\Rea\rho>1/2}}v_2(m_\rho),
 \tag{9.D1}\label{eq:sharp-dyadic-level}
\]
then
\[
 \mathsf D_k\Longleftrightarrow k\le H.
 \tag{9.D2}\label{eq:dyadic-initial-segment}
\]
Thus every dyadic realization survives through one exact finite level and all of them fail
simultaneously at \(H+1\).
\end{theorem}

\begin{proof}
Local factorization writes \(\xi(s)=(s-\rho)^{m_\rho}u(s)\) with \(u(\rho)\ne0\).
An analytic \(2^k\)-th root exists locally exactly when \(2^k\mid m_\rho\), equivalently when
the monodromy \(\omega_{\rho,k}\) is one.  Analytic continuation, the functional equation, and
Schwarz reflection transport this local criterion through the stated global and window charts.
The descended square coordinate preserves the same carrier charge on sufficiently small loops.
Finally \(\mathsf D_k\) holds exactly when \(k\le v_2(m_\rho)\) for every right-half zero,
which is \eqref{eq:dyadic-initial-segment}.
\end{proof}

This theorem explains a class of obstructions invisible to \(\nu(\rho)\).  A zero on the critical
line has unitary continuous character regardless of multiplicity, but its divided theta states,
analytic roots, and local deck torsors still depend on \(m_\rho\).  Conversely an off-line simple
zero has nonunitary continuous character but trivial ramification depth \(H=0\).  Displacement and
ramification are logically independent divisor coordinates.

\begin{theorem}[Thermal density--monodromy charge law]
\label{thm:density-monodromy}
Let \(F\) be an analytic germ at zero with values in a complex normed space, let \(q>0\), and
define
\[
 \mathfrak d_{q,F}(h)
 =\frac{\|F(e^{-qh})\|^2}{1-e^{-e^{-h}}}.
 \tag{9.M1}\label{eq:ramified-density}
\]
If \(\mathfrak d_{q,F}(h)\to\ell\) with \(0<\ell<\infty\), then \(F\) has finite positive
vanishing order \(m\) and
\[
 2qm=1.
 \tag{9.M2}\label{eq:charge-rate}
\]
If \(q=1/r\) for a positive integer \(r\), then \(r=2m\).  For a primitive \(r\)-th root
\(\zeta\), the leading deck character is therefore
\[
 \zeta^m=-1.
 \tag{9.M3}\label{eq:sign-monodromy}
\]
Consequently no finite-sheeted analytic germ which is invariant under its full deck group can
carry any positive finite first-order thermal density.  The same exclusion holds for a
nontorsion deck generator at arbitrary positive \(q\).  In particular it applies to the
four-unit normalization \(\ell=4\).
\end{theorem}

\begin{proof}
Positive limiting density excludes the identically zero germ.  Write
\(F(z)=z^mG(z)\), where \(G\) is analytic and \(G(0)\ne0\).  Since
\(1-e^{-e^{-h}}\sim e^{-h}\),
\[
 \mathfrak d_{q,F}(h)
 =e^{(1-2qm)h}\bigl(\|G(e^{-qh})\|^2+o(1)\bigr).
\]
A positive finite limit forces the exponential rate to vanish, proving
\eqref{eq:charge-rate}.  When \(q=1/r\), this gives \(r=2m\).  A primitive
\(r\)-th root satisfies \((\zeta^m)^2=1\); it cannot satisfy \(\zeta^m=1\), since that would
imply \(r\mid m\).  Hence \(\zeta^m=-1\).  Full deck invariance forces the leading coefficient
into the trivial character, contradicting \eqref{eq:sign-monodromy}.  For a nontorsion
generator, no positive power is trivial, giving the same contradiction.
\end{proof}

This theorem couples the two coordinates which the basic ramified character keeps separate.
Continuous first-order density fixes a discrete half-order, and the only compatible finite deck
sector is the sign representation.  Full descent demands the trivial representation and is
therefore impossible.  The four-unit Ward boundary is not merely another positivity test: at a
ramified analytic boundary it selects nontrivial monodromy.

\subsection{The continuous spectrum, not merely its zero endpoint}

For each divisor copy put \(z_\rho=-i(\rho-\tfrac12)\), reflect \(z_\rho\) into the closed right
half-plane, and call the result \(\lambda_\rho\).  The canonical Poisson trace is
\[
 \Theta_\xi(\tau)=\sum_\rho e^{-\tau\lambda_\rho},
 \qquad \Rea\tau>0,
 \tag{9.A1}\label{eq:poisson-trace}
\]
where zeros occur with analytic multiplicity.  Define
\[
 R_\xi=\sup_{\xi(\rho)=0}\left|\Rea\rho-\tfrac12\right|.
 \tag{9.A2}\label{eq:character-radius}
\]

\begin{theorem}[Character support--indicator spectrum]\label{thm:aperture-spectrum}
For every \(\alpha\ge0\), the following are equivalent:
\begin{enumerate}[label=\textup{(\roman*)}]
\item every Xi zero satisfies
\(\left|\Rea\rho-\tfrac12\right|\le\alpha\);
\item on some strict-right vertical line there are \(\sigma,C>0\) such that
\[
 |\Theta_\xi(\sigma+ix)|\le C e^{\alpha|x|}
 \qquad(x\in\R);
 \tag{9.A3}\label{eq:poisson-type}
\]
\item \(R_\xi\le\alpha\);
\item the supremum of the principal Hurwitz phases
\(|\arg\chi_\rho(-2\pi i)|/(2\pi)\) is at most \(\alpha\).
\end{enumerate}
Consequently the common admissible aperture spectrum is exactly
\[
 \{\alpha\ge0:\textup{(i)--(iv) hold}\}=[R_\xi,\infty).
 \tag{9.A4}\label{eq:aperture-ray}
\]
The same half-plane condition is the common parameter in the Hedenmalm Gaussian Hankel,
inverse-tilted prime-source, critical-Mellin, and prime--Gamma negative-energy charts.
\end{theorem}

\begin{proof}
The imaginary part of \(\lambda_\rho\) has absolute value
\(|\Rea\rho-\tfrac12|\).  The support--indicator theorem for an absolutely convergent discrete
Poisson trace therefore identifies \eqref{eq:poisson-type} with the uniform bound in (i).
Condition (iii) is the definition of the least such uniform bound.  On the open critical strip,
the principal-period identity
\[
 \frac{|\arg\chi_\rho(-2\pi i)|}{2\pi}
 =\left|\Rea\rho-\tfrac12\right|
\]
has no phase aliasing, proving (iv).
\end{proof}

This is stronger than a decorative equivalence at \(\alpha=0\).  It determines every horizontal
zero-free chamber and proves that all admissible apertures form one closed upper ray.  RH is only
the degenerate endpoint \(R_\xi=0\).

\begin{corollary}[Mass-graded aperture distortion]
\label{cor:mass-distortion}
Let \(a\ge\tfrac12\) and let the mass-\(a\) transform replace each Poisson point
\(\lambda_\rho\) by the principal root
\(\lambda_{\rho,a}=\sqrt{\lambda_\rho^2+a^2}\), so that
\[
 \Theta_{\xi,a}(\tau)=\sum_\rho e^{-\tau\lambda_{\rho,a}}.
\]
If every Xi zero has horizontal displacement at most \(\lambda\ge0\), then
\(\Theta_{\xi,a}\) has vertical exponential type at most \(\lambda\) on a strict-right line.
Conversely, if that mass trace has type at most \(\lambda\), then every Xi zero satisfies
\[
 \left|\Rea\rho-\tfrac12\right|
 \le \sqrt{\frac{9+a^2}{9}}\,\lambda.
 \tag{9.A5}\label{eq:mass-distortion}
\]
\end{corollary}

\begin{proof}
Mass shifting contracts the imaginary coordinate of every centered spectral point, so the
original width bound gives the same trace-type bound.  Faithful inverse mass transfer obeys
\[
 9\left|\Rea\rho-\tfrac12\right|^2
 \le(9+a^2)\lambda^2,
\]
which is equivalent to \eqref{eq:mass-distortion}.
\end{proof}

Thus exact conservative charts preserve the aperture ray, whereas a faithful nonconservative
consumer returns it with an explicit condition number.  This is the quantitative data-processing
form of the character law: information need not disappear, but its metric can be distorted.

\begin{theorem}[Bifiltered ramified-character boundary law]
\label{thm:singular-equalizer}
The Riemann source has a four-stage boundary architecture.
\begin{enumerate}[label=\textup{(\roman*)}]
\item \emph{Bifiltered singular acquisition.}  Valuation depth and prime-normal breadth admit no
faithful additive scalarization.  Every nonzero positive momentum-symmetric theta-divisor
detector is a singular atomic rigging and is nonclosable through ordinary \(L^2\); no entire
heat-node symbol isolates one source atom.
\item \emph{Odd orientation jet.}  The unique positive Julia lift restores the prime logarithmic
normal derivative erased by the scalar Gamma contraction.  Its finite packets converge to the
maximally singular prime clock, and the pointed Fisher--Ward trace is odd in this unsquared
tangent and pairs it exactly with cutoff curvature.  Squaring retains the full logarithmic
current but loses this odd lift; no finite-dimensional reporter or finite local polynomial frame
recovers it.
\item \emph{Ramified divisor output.}  The divisor output is
\((m_\rho,\chi_\rho)\).  Global neutrality of the continuous inner
duality--adjunction fracture is equivalent simultaneously to independent-quadrant
Wold--Poisson boundedness on every divisor fibre, compactness of every seam-weighted
Sonin/Calkin chart, and positivity of every fixed-edge four-unit Ward Gram.
Multiplicity is realized at every dyadic level by analytic
root descent, local loop monodromy, and window deck transport, with the exact sharp level in
Theorem~\ref{thm:dyadic-descent}.  Positive finite first-order thermal density couples the axes
by \(2qm=1\); on a finite \(r\)-sheeted branch it selects sign monodromy and excludes full deck
descent.
\item \emph{Consumer compression.}  The trace-class energy \(\mathcal E_\chi\) is a faithful
quadratic detector of the global unitary equalizer but is reflection-even.  Every strict-right
zero has one positive rate \(\nu=2c\) which simultaneously produces
Wold--Poisson quadrant blow-up, prime-seam incoherence, a noncompact Sonin/Calkin corner, the
one-node Ward failure \eqref{eq:ward-local-failure}, and the finite-versus-flat Rees fracture of
Theorem~\ref{thm:trichotomy}.
Every faithful norm-conservative equivariant descendant forces that continuous rate to
vanish, while the complete quantitative phase is the aperture ray
\([R_\xi,\infty)\), not only its zero endpoint.  Faithful mass deformation preserves this
information with the explicit distortion factor \(\sqrt{(9+a^2)/9}\).
\end{enumerate}
Therefore no one scalar regular ordinary-Hilbert completion in the stated source-core and
consumer classes can acquire the Xi divisor, retain its prime-normal source grade and
multiplicity, preserve the odd orientation jet, and conservatively reflect its continuous
character.  The exact surviving architecture is a bifiltered singular source, an odd
orientation jet, and the ramified character module
\(\{(m_\rho,\chi_\rho)\}_\rho\), followed only then by quadratic or regular consumers.
\end{theorem}

\begin{proof}
Part (i) combines Theorems~\ref{thm:bifilter}, \ref{thm:theta-rigging}, and
\ref{thm:heat-node}.  Part (ii) is
Theorems~\ref{thm:odd-orientation} and \ref{thm:parity-splitting}, together with
Corollary~\ref{cor:finite-reader}.
The Wold equivalence in part (iii) is
Theorem~\ref{thm:wold-poisson}; the Calkin equivalence is
Theorem~\ref{thm:calkin-equalizer}; the Ward equivalence is
Theorem~\ref{thm:ward-equalizer}.  Their equalizer is the common identity
\(\Delta(\rho)=\nu(\rho)=2(\Rea\rho-\tfrac12)\); its ramification statements are
Theorems~\ref{thm:dyadic-descent} and \ref{thm:density-monodromy}.  Part (iv) combines
Theorem~\ref{thm:character-energy}, the strict versions of the continuous consumer theorems,
Corollary~\ref{cor:xi-superselection}, and
Theorems~\ref{thm:conservative-descent} and \ref{thm:aperture-spectrum}, with
Corollary~\ref{cor:mass-distortion}.
\end{proof}

\section{Interpretation for RH}\label{sec:interpretation}

Theorems~\ref{thm:inner-outer}, \ref{thm:duality-fracture},
\ref{thm:critical-resolution}, \ref{thm:ht-creation}, \ref{thm:graph-weld},
\ref{thm:trichotomy}, \ref{thm:essential}, and \ref{thm:singular-equalizer} identify a common
mechanism behind several earlier obstructions.

\begin{center}
\begin{tabular}{@{}lll@{}}
\toprule
Source operation & retained datum & discarded datum\\
\midrule
fixed exponent-depth limit & \(\zeta\) bulk & reciprocal divisor\\
inverse-depth normal cone & \(1/(2\xi)\) & positive coefficient chart\\
neutral same-prime fusion & exact coefficient law & finite differential closure\\
additive source totalization & one mixed degree & depth--breadth pairing\\
Gamma scalar contraction & quadratic contact & prime-log first jet\\
positive Julia lift & oriented prime-log jet & bounded global recovery\\
pointed Fisher--Ward trace & signed cutoff curvature & quadratic energy\\
squared completed carrier & divisor/log current & odd reflection lift\\
finite jet reporter & selected derivative packet & complete current tower\\
score/overlap graph & Xi outer connection & bounded positive completion\\
two-endpoint normal blow-up & inner pole orientation & global normalization\\
centered Euler contraction & affine residue & uniform Herglotz sign\\
continuous modular character & full spectral point & divisor multiplicity\\
dyadic deck character & \(2\)-adic multiplicity & horizontal displacement\\
thermal boundary density & half-order sign sector & full deck descent\\
prime-character transport & endpoint separation & nontrivial loop holonomy\\
even scalarization & divisor magnitude & reflection orientation\\
Xi--Hurwitz energy & global unitary equalizer & side of each defect\\
Kakutani quotient & singularity & finite versus flat face\\
Gamma sewing & every finite prime packet & bounded Hilbertization\\
finite-Rees transport & each Rees grade & all-orders type crossing\\
comoving clock restriction & finite traveling profile & independent-quadrant rate\\
ordinary \(L^2\) closure & cyclic theta source & nonzero divisor rigging\\
entire heat-node symbol & boundary accumulation & single-node separation\\
Sonin Calkin quotient & coherent prime seams & off-equalizer compactness\\
Ward Schur complement & four-unit boundary ball & strict-right character\\
norm-conservative descent & unitary character & nonzero real rate\\
mass deformation & conserved charge & exact aperture metric\\
\bottomrule
\end{tabular}
\end{center}

The unifying statement is:

\begin{quote}
\emph{The complete pre-positivity datum is not a scalar zero statistic.  It is a bifiltered
singular source with an odd orientation jet, whose divisor output is the ramified character
\((m_\rho,\chi_\rho)\).  Source depth and prime-normal breadth cannot be collapsed to one additive
degree.  The Julia normal derivative and pointed Fisher--Ward transgression retain the
unsquared prime-log direction; its even square retains the logarithmic current but not the
reflection lift.  Singular theta and squarefree graphs acquire the divisor.  The continuous character
\(\chi_\rho\) retains spectral position; its real rate controls Wold--Poisson boundedness,
Sonin/Calkin compactness, Ward positivity, Rees type, and the complete aperture ray.  The discrete
coordinate \(m_\rho\) controls analytic roots and local deck descent.  Their first coupling is the
inner--outer residue \(-m_\rho(\rho-\tfrac12)\); positive thermal boundary density couples them
again through \(2qm=1\) and sign monodromy.  Every faithful conservative descendant forces
the continuous character to be unitary.  Quadratic character energy detects this equalizer but
erases its side.  The former no-go theorems are the different quotient shadows obtained by
forgetting source breadth, odd orientation, singular acquisition, continuous rate, or
ramification.}
\end{quote}

This statement answers several no-go results without turning them into slogans.  The balanced
Euler contraction cannot be Herglotz because it retains the additive fracture and therefore has
opposite residues at reflected real zeros; the uncontracted logarithmic one-form is the
sign-compatible object.  Prime translation is a faithful detector but a flat coboundary, because
exponentials of an additive character have trivial closed-orbit holonomy; detection alone cannot
force the character to vanish.  A uniformly bounded positive Gram completion cannot carry the
score through a nonzero divisor; the graph crosses that boundary by the explicitly singular
division whose inverse-square cost is \eqref{eq:gram-blowup}.  Finally, every scalar observer
depending only on \(a^2\) or \(|a|\) is invariant under \(a\mapsto-a\), so it factors through an
even quotient of the duality--adjunction fracture.  Theorem~\ref{thm:character-energy} proves
that this is not a defect of one positivity construction: the canonical global positive
compression is complete for the equalizer and still reflection-even.

Theorem~\ref{thm:singular-equalizer} also explains the obstructions that this source-completeness
statement previously left outside.  Wold--Poisson ancestry is a two-clock representation of the
same additive fracture; locking the clocks to a characteristic makes every rate look finite.
Sonin/Calkin compactness is the quotient norm of prime-seam incoherence, not an independent
spectral condition.  The four-unit Ward law is the final Schur consumer: a strict-right rate
produces a one-node negative diagonal.  Hedenmalm's dense radical lives at acquisition:
even a neutral real-divisor sampler cannot descend to ordinary \(L^2\), and heat-node
accumulation prevents an analytic delta reader.  Hedenmalm multiplicity jets and the Xi
root/deck ladder live instead on the ramification coordinate; horizontal unitarity cannot see
them.  The rank-one co-Poisson conormal supplies a scalar seam but cannot generate a
prime-labelled Calkin frame because scalarization has erased source breadth.  Hence neither
adding more finite depth coordinates nor changing only the final positivity test repairs the
architecture.

The Julia and pointed-Fisher clauses identify what those negative results were missing.
The Gamma contraction reaches the lossless boundary only to second order; its positive
two-channel completion restores the first prime-log jet.  That jet is necessarily singular in
the global Hardy topology, but its pointed Ward transgression remains odd.  Orientation is
therefore neither absent nor supplied by the final four-unit Gram: it exists one categorical
level earlier, before quadratic scalarization.  The Ward Gram is a consumer of the character,
not its source.

The parity theorem explains two further no-go families.  Squared completed Xi can reconstruct
the divisor and logarithmic current because evenization does not destroy the logarithmic
differential field.  What it destroys is the odd lift, and every differentiable even reader has
zero first jet.  A fixed finite reporter panel therefore misses a nonzero derivative packet,
while a finite constant-coefficient local frame would force an impossible real-shift recurrence
of the Xi divisor.  These are not failures to detect zeros; they are failures to retain the
oriented lift of a current already detected.

The density--monodromy law also shows why finite ramification does not repair four-unit
positivity.  Positive first-order density selects the nontrivial sign representation, whereas
full deck descent requires the trivial representation.  Conversely, mass deformation is not a
no-go: it is a faithful but nonisometric consumer whose exact condition number measures the loss.
The architecture therefore distinguishes erasure, singular preservation, nonconservative
distortion, and outright representation incompatibility.

Large terminal equivalence packages now have a precise status as well.  A single positive Sonin
collar, trace/Hankel positivity, continued-Cauchy passivity, geometric phase certificates, and
the other terminal charts all collapse because they test the zero locus of the same quadratic
character energy.  Their simultaneous equivalence is real, but it creates no additional source
coordinate and no mechanism forcing zero energy.  The object theorem therefore explains those
packages without using their size as a novelty claim.

The divisor class enters here through two explicitly singular boundary operations: the
nonclosable theta rigging and score/overlap division.  Once acquired, its character cannot be
altered by a faithful conservative intertwiner.  Acquisition is singular; descent is
conservative only on the unitary equalizer.

The result gives a closed categorical restriction on arguments implemented through these
source cores and reflected valuation sectors.  The complete architecture is
\[
 \begin{array}{c}
 \text{depth--breadth prime source}
 \longrightarrow\text{countably free bifiltered fusion tower}\\
\hspace{8mm}\downarrow\ \text{positive Julia lift / odd Fisher--Ward jet}\\
 \hspace{10mm}\searrow\ \text{singular score/overlap graph}
 \longrightarrow\text{Xi outer connection}\\[-1mm]
 \text{theta source}\longrightarrow\text{singular real-divisor rigging}
 \hspace{8mm}\nearrow\\[1mm]
 \text{ramified divisor module }\{(m_\rho,\chi_\rho)\}_\rho\\
 \swarrow\ \begin{array}{c}
 \text{multiplicity/deck tower}\\
 \text{analytic roots and local loops}
 \end{array}
 \qquad
 \searrow\ \begin{array}{c}
 \text{continuous aperture ray}\\
 \text{Wold--Poisson, Sonin/Calkin}\\
 \text{Ward, energy, Fisher--Gamma/Rees}
 \end{array}
 \end{array}
 \tag{10.1}\label{eq:architecture}
\]
Theorem~\ref{thm:main} ties every arrow in \eqref{eq:architecture} to the same exponent source and
identifies the complete divisor output in inner, outer, clock, quotient, character, and deck
charts.  The source bifiltration preserves prime-normal breadth, while the theta leg proves why
divisor acquisition must be a singular rigging rather than another bounded completion.
Corollary~\ref{cor:xi-superselection} proves the exhaustive conclusion for the stated transport
class: no bounded finite-Rees completion can perform the reflection step.

\section{Conclusion}

One source law now explains the previously separate bulk, divisor, monodromy, shell, window, and
transport observations.  Positive bounded-exponent Euler sources have three distinct scales:
their fixed-scale limit is \(\zeta\), their inverse-depth limit is reciprocal zeta, and their
depth derivative is the Xi logarithmic current after completion.  The centered residue of that
current is multiplicity times the full affine displacement from the critical line.

The missing transition has now been identified exactly.  A squarefree Fisher score divided by
its own overlap, with the forced doubled-zeta correction, reconstructs
\(-2\xi'/\xi\).  Balanced Euler contraction gives the outer current, while two-endpoint moving-zero
blow-up gives the inner rational profile.  Their transition coefficient is the same divisor
charge.  Source inverse and analytic adjoint differ by twice the horizontal displacement, and
every prime seam is the exponential image of those same inner poles.

The same displacement is the logarithmic modulus of the cofinal character of a positive
hypertranscendental shell.  Prime two and prime three recover the entire complex coordinate.
Compact prime-window dilation and Hurwitz monodromy recover the same charge independently.
These are exact identities between source constructions, not zero-set definitions.

The depth-one source also identifies its own completion.  Its positive valuation atoms are
Fisher-damped Gamma coordinates.  Removing that damping restores prime-harmonic divergence;
canonical Gamma sewing preserves summability exactly, while its boundary symbols have
arbitrarily large weighted-Vandermonde rank.  Gamma is therefore an infinite Rees resolution of
the source, not an omitted bounded coordinate.

No finite substitute has the same closure properties.  Same-prime neutral fusion forces a
countably free hypertranscendental valuation--residue tower.  Divisor differentiation lowers
overlap order and forces inverse-square positive-Gram blow-up.  Reflection sends one sector into
an asymptotic species inaccessible to every fixed finite Rees order.  These are exact theorems in
three different categories, anchored to one ordinary-prime source by the squarefree graph.

The source also supplies the orientation that its even quotients conceal.  The Gamma scalar
contraction approaches its lossless edge quadratically; the positive Julia lift restores the
first normal jet \(\log p\), and the pointed Fisher--Ward trace is odd in that jet.  Globally,
the prime-clock graph is maximally singular and every bounded recovery diverges.  After divisor
formation, the Xi--Hurwitz character energy is a finite positive compression with exactly the
same unitary zero locus, but reflection invariance removes the sign.  This separates three
questions that earlier criteria conflated: producing orientation, detecting its failure, and
forcing its vanishing.  The even-current/odd-orientation theorem makes the separation exact:
squaring preserves hypertranscendence and the logarithmic-current field, while every even
differentiable reader loses the first orientation jet.  Finite reporter and finite local-frame
no-gos are consequences of that split, not evidence that the divisor itself is inaccessible.

Reflection then exposes a genuine type fracture.  At every hypothetical right-half Xi zero, the
corresponding prime-valuation sector is flat beyond all finite Rees orders, while its reflected
partner is half-normal.  The flat radical is invariant under arbitrary-dimensional bounded
finite-Rees transport, and exact sewing has essential operator norm.

The broader law is an oriented bifiltered ramified-character law.  Before the divisor, valuation
depth and prime-normal breadth are independent and no one additive total degree is faithful.
Their odd Julia/Fisher jet carries the prime-log direction before scalar squaring.  At the
divisor, the complete output is
\[
 (m_\rho,\chi_\rho),\qquad
 \chi_\rho(z)=e^{z(\rho-1/2)}.
\]
The continuous rate \(\nu(\rho)=2(\Rea\rho-\tfrac12)\) is simultaneously the inner pole fracture,
independent-clock rate, prime-seam Calkin coefficient, and obstruction seen by the fixed-edge
Ward Schur complement.  A nonzero rate does not produce four unrelated failures: it produces one
nonunitary character in four regular consumer categories.  Its full quantitative phase is the
closed aperture ray \([R_\xi,\infty)\).

Vertical divisor acquisition behaves differently.  A nonzero positive theta-null form is forced
onto the real Xi divisor, but the theta source is cyclic in ordinary \(L^2\).  Its radical is
therefore dense and its form nonclosable; accumulating heat nodes also forbid an analytic delta
reader.  Singular divisor rigging is necessary even when horizontal rate vanishes.  The
independent multiplicity \(m_\rho\) then controls the full dyadic root/deck phase and its exact
sharp breakdown level.  Divisor detection is not a bounded precursor of reflection positivity,
and reflection positivity is not a way to create or ramify the divisor.

Continuous density and discrete ramification meet in a second universal charge relation.
Positive finite thermal density at rate \(q\) forces \(2qm=1\); finite-sheeted realizations
therefore live in the sign deck sector and cannot descend through the full deck group.  By
contrast, mass-shifted Poisson traces remain faithful: their exponential type returns the
original horizontal width with the explicit condition number \(\sqrt{(9+a^2)/9}\).

The Riemann hypothesis remains the assertion that every nontrivial divisor charge is neutral.
The contribution here is the prior classification: source bidegree, singular acquisition,
odd orientation jet, continuous character, discrete ramification, their
multiplicity-weighted inner--outer coupling, parity split, thermal charge relation, quadratic
detector, and universal conservative-descent equalizer.  Within the stated source classes, the correct object is not one
Hilbert completion but a bifiltered singular source, its odd normal jet, and its ramified
character module, together with countably free hypertranscendental and infinite-Rees resolutions.

\begin{thebibliography}{99}

\bibitem{ArratiaBarbourTavare2003}
R.~Arratia, A.~D.~Barbour, and S.~Tavar\'e,
\newblock \emph{Logarithmic Combinatorial Structures: A Probabilistic Approach},
\newblock EMS Monographs in Mathematics, European Mathematical Society, 2003.

\bibitem{Burnol2004}
J.-F.~Burnol,
\newblock Two complete and minimal systems associated with the zeros of the Riemann zeta function,
\newblock \emph{J. Th\'eor. Nombres Bordeaux} \textbf{16} (2004), 65--94.
\newblock \url{https://doi.org/10.5802/jtnb.434}.

\bibitem{Connes1999}
A.~Connes,
\newblock Trace formula in noncommutative geometry and the zeros of the Riemann zeta function,
\newblock \emph{Selecta Math. (N.S.)} \textbf{5} (1999), 29--106.
\newblock \url{https://doi.org/10.1007/s000290050042}.

\bibitem{ConnesConsani2021}
A.~Connes and C.~Consani,
\newblock The scaling Hamiltonian,
\newblock \emph{J. Operator Theory} \textbf{85} (2021), 257--276.
\newblock \url{https://arxiv.org/abs/1910.14368}.

\bibitem{Connes2026}
A.~Connes,
\newblock The Riemann hypothesis: past, present and a letter through time,
\newblock arXiv:2602.04022 (2026).
\newblock \url{https://arxiv.org/abs/2602.04022}.

\bibitem{GauthierXarles2009}
P.~M. Gauthier and X.~Xarles,
\newblock Perturbations of \(L\)-functions with or without non-trivial zeros off the critical
line,
\newblock arXiv:0911.5135 (2009).
\newblock \url{https://arxiv.org/abs/0911.5135}.

\bibitem{Hedenmalm2026}
H.~Hedenmalm,
\newblock Spectral interpretation of Riemann zeta zeros,
\newblock arXiv:2606.17494 (2026).
\newblock \url{https://arxiv.org/abs/2606.17494}.

\bibitem{Hellinger1909}
E.~Hellinger,
\newblock Neue Begr\"undung der Theorie quadratischer Formen von unendlichvielen
Ver\"anderlichen,
\newblock \emph{J. Reine Angew. Math.} \textbf{136} (1909), 210--271.

\bibitem{Kakutani1948}
S.~Kakutani,
\newblock On equivalence of infinite product measures,
\newblock \emph{Ann. of Math. (2)} \textbf{49} (1948), 214--224.
\newblock \url{https://doi.org/10.2307/1969123}.

\bibitem{Meyer2005}
R.~Meyer,
\newblock A spectral interpretation for the zeros of the Riemann zeta function,
\newblock in \emph{Seminars Winter Term 2004/2005}, Universit\"atsdrucke G\"ottingen,
2005, pp.~117--137.
\newblock \url{https://arxiv.org/abs/math/0412277}.

\bibitem{Nagoshi2016}
H.~Nagoshi,
\newblock The hypertranscendence of \(L\)-functions for
\(\mathrm{GL}_m(\mathbb A_{\mathbb Q})\),
\newblock \emph{Bull. Aust. Math. Soc.} \textbf{93} (2016), 388--399.
\newblock \url{https://doi.org/10.1017/S000497271500129X}.

\bibitem{Ostrowski1920}
A.~Ostrowski,
\newblock \"Uber Dirichletsche Reihen und algebraische Differentialgleichungen,
\newblock \emph{Math. Z.} \textbf{8} (1920), 241--298.
\newblock \url{https://doi.org/10.1007/BF01206523}.

\bibitem{Pong2016}
W.-Y.~Pong,
\newblock Applications of differential algebra to algebraic independence of arithmetic
functions,
\newblock \emph{Acta Arith.} \textbf{172} (2016), 149--173.
\newblock \url{https://doi.org/10.4064/aa8112-12-2015}.

\bibitem{Postnikov1949}
A.~G.~Postnikov,
\newblock On the differential independence of Dirichlet series,
\newblock \emph{Dokl. Akad. Nauk SSSR} \textbf{66} (1949), 561--564
\newblock (in Russian).

\bibitem{ShapiroSparer1986}
H.~N.~Shapiro and G.~H.~Sparer,
\newblock On algebraic independence of Dirichlet series,
\newblock \emph{Comm. Pure Appl. Math.} \textbf{39} (1986), 695--745.
\newblock \url{https://doi.org/10.1002/cpa.3160390602}.

\bibitem{Schwartz1966}
L.~Schwartz,
\newblock \emph{Th\'eorie des distributions},
\newblock Hermann, Paris, 1966.

\bibitem{Titchmarsh1986}
E.~C.~Titchmarsh,
\newblock \emph{The Theory of the Riemann Zeta-Function}, 2nd ed.,
\newblock revised by D.~R.~Heath-Brown, Oxford University Press, 1986.

\bibitem{Wiener1932}
N.~Wiener,
\newblock Tauberian theorems,
\newblock \emph{Ann. of Math. (2)} \textbf{33} (1932), 1--100.

\end{thebibliography}

\end{document}
